\documentclass[11pt,oneside]{article}

\usepackage[a4paper,margin=27mm,headheight=15pt]{geometry}
\usepackage[T1]{fontenc}
\usepackage[utf8]{inputenc}
\IfFileExists{libertinus.sty}{\usepackage{libertinus}}{\usepackage{lmodern}}
\usepackage{microtype}
\usepackage{amsmath,amssymb,amsthm,mathtools,mathrsfs}
\usepackage{bm}
\usepackage{booktabs,longtable,array,tabularx}
\usepackage{graphicx}
\usepackage{pgfplots}
\usepackage{xcolor}
\usepackage{enumitem}
\usepackage{fancyhdr}
\usepackage{titlesec}
\usepackage{listings}
\usepackage{xurl}
\usepackage{hyperref}
\usepackage[nameinlink,noabbrev]{cleveref}

\definecolor{linkblue}{RGB}{25,75,135}
\definecolor{shadegray}{RGB}{246,247,249}
\definecolor{rulegray}{RGB}{195,201,208}
\pgfplotsset{compat=1.18}
\hypersetup{
  colorlinks=true,
  linkcolor=linkblue,
  citecolor=linkblue,
  urlcolor=linkblue,
  pdftitle={Computability of the Inverse Fabius Function},
  pdfauthor={Research report prepared for Vladimir Reshetnikov's ProveIt project},
  pdfsubject={Sequential computability, effective uniform continuity, exact inverse moduli, Thue--Morse splines, and certified inversion},
  pdfkeywords={Fabius function, inverse Fabius function, computable analysis, Grzegorczyk computability, Thue--Morse, Rvachev up function, q-binomial, sinc product, Lambert W}
}

\pagestyle{fancy}
\fancyhf{}
\fancyhead[L]{\small Computability of the inverse Fabius function}
\fancyhead[R]{\small\nouppercase{\leftmark}}
\fancyfoot[C]{\thepage}
\renewcommand{\headrulewidth}{0.4pt}

\titleformat{\section}{\Large\bfseries}{\thesection}{0.7em}{}
\titleformat{\subsection}{\large\bfseries}{\thesubsection}{0.7em}{}
\titleformat{\subsubsection}{\normalsize\bfseries}{\thesubsubsection}{0.7em}{}
\setcounter{secnumdepth}{3}
\setcounter{tocdepth}{2}
\setlist[itemize]{leftmargin=2em,itemsep=0.25em,topsep=0.4em}
\setlist[enumerate]{leftmargin=2.3em,itemsep=0.25em,topsep=0.4em}

% Long unbreakable inline mathematics occasionally leaves a line with no legal
% break point.  A small emergency stretch lets TeX loosen such a line rather
% than let it protrude into the margin; it applies only to lines that would
% otherwise be overfull.
\setlength{\emergencystretch}{2em}

\newtheorem{theorem}{Theorem}[section]
\newtheorem{proposition}[theorem]{Proposition}
\newtheorem{lemma}[theorem]{Lemma}
\newtheorem{corollary}[theorem]{Corollary}
\newtheorem{conjecture}[theorem]{Conjecture}
\theoremstyle{definition}
\newtheorem{definition}[theorem]{Definition}
\newtheorem{algorithm}[theorem]{Algorithm}
\newtheorem{example}[theorem]{Example}
\theoremstyle{remark}
\newtheorem{remark}[theorem]{Remark}
\newtheorem{warning}[theorem]{Warning}

\newcommand{\R}{\mathbb R}
\newcommand{\C}{\mathbb C}
\newcommand{\Q}{\mathbb Q}
\newcommand{\N}{\mathbb N}
\newcommand{\Z}{\mathbb Z}
\newcommand{\Up}{\operatorname{up}}
\newcommand{\sinc}{\operatorname{sinc}}
\newcommand{\wt}{\operatorname{wt}_2}
\newcommand{\e}{\mathrm e}
\newcommand{\ii}{\mathrm i}
\newcommand{\dd}{\,\mathrm d}
\newcommand{\Li}{\operatorname{Li}}
\newcommand{\Var}{\operatorname{Var}}
\newcommand{\E}{\mathbb E}
\newcommand{\Prob}{\mathbb P}
\newcommand{\bigO}{\mathcal O}
\newcommand{\smallo}{o}
\newcommand{\supp}{\operatorname{supp}}
\newcommand{\sgn}{\operatorname{sgn}}
\newcommand{\lcm}{\operatorname{lcm}}
\newcommand{\vtwo}{v_2}
\newcommand{\qbinom}[3]{\genfrac{[}{]}{0pt}{}{#1}{#2}_{#3}}
\newcommand{\repo}{\nolinkurl{Analysis/FabiusFunction}}
\newcommand{\lean}[1]{%
  \ifmmode\text{\nolinkurl{#1}}\else\nolinkurl{#1}\fi}
\newcommand{\leanevaltwo}[1]{%
  \nolinkurl{integral_eval}\texttt{\textsubscript{2}}\nolinkurl{#1}}

\newenvironment{proofidea}{\par\smallskip\noindent\textit{Proof architecture.}\ }{\hfill$\square$\par\smallskip}
\newenvironment{boxedremark}{%
  \par\smallskip\noindent\begin{center}\begin{minipage}{0.94\textwidth}%
  \color{black}\hrule\vspace{0.6em}\small
}{%
  \vspace{0.6em}\hrule\end{minipage}\end{center}\smallskip
}

\lstdefinestyle{pseudo}{
  basicstyle=\ttfamily\small,
  backgroundcolor=\color{shadegray},
  frame=single,
  rulecolor=\color{rulegray},
  columns=fullflexible,
  keepspaces=true,
  showstringspaces=false,
  xleftmargin=0.7em,
  xrightmargin=0.7em,
  aboveskip=0.8em,
  belowskip=0.8em
}

% Report-local layout, notation, environments, and listing extension.
\allowdisplaybreaks[2]
\numberwithin{equation}{section}

\newtheorem{problem}[theorem]{Research problem}

\crefname{problem}{research problem}{research problems}
\Crefname{problem}{Research problem}{Research problems}

\newcommand{\abs}[1]{\left\lvert #1\right\rvert}
\newcommand{\norm}[1]{\left\lVert #1\right\rVert}
\newcommand{\set}[1]{\left\{#1\right\}}
\newcommand{\floor}[1]{\left\lfloor #1\right\rfloor}
\newcommand{\ceil}[1]{\left\lceil #1\right\rceil}
\newcommand{\pospart}[1]{\left(#1\right)_{+}}
\newcommand{\TM}{\varepsilon}
\newcommand{\Inv}{G}
\newcommand{\InvT}{\overline G}
\newcommand{\Clamp}{\operatorname{clamp}}
\newcommand{\arrivalcommit}{\nolinkurl{40fdea4cc0a728189f357389e3f114a2cb00e561}}
\newenvironment{statusbox}[1]{%
  \par\smallskip\noindent\begin{center}\begin{minipage}{0.94\textwidth}%
  \color{black}\hrule\vspace{0.55em}\small\noindent\textbf{#1}\par\smallskip
}{%
  \vspace{0.55em}\hrule\end{minipage}\end{center}\smallskip
}
\lstdefinestyle{reportpseudo}{style=pseudo,breaklines=true}

\title{\bfseries Computability of the Inverse Fabius Function\\[0.4em]
\large Sequential computation, effective uniform continuity,\\
exact inverse moduli, and certified Thue--Morse bisection}
\author{Research report prepared for Vladimir Reshetnikov's \texttt{ProveIt} project}
\date{30 August 2026\\[0.25em]
\small Arrival snapshot: \arrivalcommit; current-corpus status audited on intake}

\begin{document}
\maketitle

\begin{abstract}
Let $F$ be the canonical bounded Fabius function, viewed as the strictly
increasing distribution function on $[0,1]$, and let $G=F^{-1}$.  At the
paper level, we prove constructively that $G$ is a computable real function in the
Grzegorczyk/Wikipedia sense \cite{G1957,Weihrauch2000,WikipediaCRF}: it maps every computable real sequence to a
computable real sequence, and it possesses a recursive modulus of uniform
continuity.  The same theorem holds for the repository's totalized inverse
$\overline G:\R\to[0,1]$, obtained by clamping the argument to $[0,1]$.

The proof is quantitative.  First, the centered finite random sums give exact
rational Thue--Morse splines $S_N$ satisfying the global certificate
$\lvert S_N-F\rvert\le 2^{-N}$.  Thus $F$ can be evaluated at a rational point
to any requested precision by a finite integer computation.  Second, the
Fabius density is symmetric and strictly unimodal.  It follows that among all
subintervals of $[0,1]$ of a fixed nonzero length $h$, exactly the endpoint
interval starts minimize the Fabius mass:
\[
  \min_{0\le x\le1-h}\bigl(F(x+h)-F(x)\bigr)=F(h).
\]
More precisely, equality holds only for $x=0$ or $x=1-h$ (these coincide when
$h=1$); for $h=0$ every admissible start is degenerate.  Consequently the
inverse-gap bound also has a complete equality classification, not merely a
sharp endpoint witness.  The exact modulus of continuity of the inverse is the
inverse itself, with the natural saturation at unit input distance:
\[
  \sup_{\substack{u,v\in[0,1]\\\lvert u-v\rvert\le\eta}}
  \lvert G(u)-G(v)\rvert=G(\min(\eta,1)).
\]
This identity reduces effective continuity to a positive lower bound for
$F(h)$.  A direct box event in the defining random series gives, for every
$r\ge0$,
\[
  F(2^{-r})\ge
  \Delta_r:=2^{-r(r+1)/2}(r+1)^{-(r+1)}.
\]
Hence $\lvert u-v\rvert<\Delta_r$ implies
$\lvert\overline G(u)-\overline G(v)\rvert<2^{-r}$, yielding an explicit
primitive-recursive modulus.  Combined with the standard flatness upper bound,
this also determines the optimal binary input precision for $p$ output bits:
\[
  \mu(p)=\frac{p^2}{2}+p\log_2p+\bigO(p).
\]
Finally, a tolerant bisection algorithm uses certified approximations to both
$F(c)$ and the target value.  It never asks an undecidable exact-equality
question: when comparison is inconclusive, the inverse modulus certifies that
the midpoint is already sufficiently accurate.  This supplies a uniform
name-transformer and proves sequential computability at the
mathematical-report level.

The report audits the recursive LaTeX corpus under
\path{Analysis/FabiusFunction/docs}, separates imported
repository results from the arguments proved here, relates the construction to
$q$-binomial and sinc-product representations, Lambert-$W$ inverse asymptotics,
Legendre expansions, and Bell/Bernoulli coefficient calculi.  It records both
the completed post-publication Lean crosswalks for the structural inverse
modulus and effective uniform continuity, and the remaining roadmap for a
computable name-transformer.  An
exact-rational Python supplement reproduces the finite-spline and modulus
calculations.
\end{abstract}

\begin{statusbox}{Status, snapshot, and formalization boundary}
Every statement labelled theorem, proposition, lemma, or corollary below is
proved in this report or explicitly imported with a source pointer.  The
repository audit and all relative-novelty statements remain pinned to commit
\arrivalcommit.  At that snapshot, the main computability theorem, the least
fixed-length Fabius-increment inequality, the exact inverse modulus, and the
closed elementary modulus were absent under the formulations used here.

A post-publication Lean development now formalizes three layers.  The structural
inverse-modulus layer lives in
\path{Analysis/FabiusFunction/Lean/FabiusFunction/InverseModulus.lean}.  It
proves the global weak shape and maximal strict branches of fixed-length
Fabius increments, the endpoint lower bound and complete equality locus,
constrained forward superadditivity with its equality cases, local and global
inverse-gap inequalities, unit-interval inverse-gap rigidity, global inverse
subadditivity, attained exact moduli on both $[0,1]$ and $\R$, and the exact
effective-injectivity threshold.  The downstream module
\path{Analysis/FabiusFunction/Lean/FabiusFunction/FabiusInverseEffectiveContinuity.lean}
proves a one-term recurrence lower bound stronger than the elementary
$\Delta_r$ estimate, derives the exact numerical inequality
$\Delta_r\le F(2^{-r})$ and the corresponding strict and closed dyadic inverse
moduli, and packages a primitive-recursive factorial-denominator witness for
\texttt{EffectivelyUniformContinuous}.  The further refinement
\path{Analysis/FabiusFunction/Lean/FabiusFunction/FabiusInverseLogarithmicModulus.lean}
formalizes the least logarithmic order $r(n)$, composes both denominator
families with it, and packages the report-exact $d(n)$ as well as the stronger
factorial alternative.  Exhaustive declaration-level crosswalks are given in
\cref{sec:lean-status}.

The probabilistic box-event \emph{proof} of the $\Delta_r$ inequality is not
formalized: Lean reaches the same inequality by the stronger recurrence
route.  The exact ceiling denominator $d_*$ and its qualified fixed-target
minimality, tolerant bisection, sequential computability, the combined
computable-real-function theorem, and the input-precision asymptotics remain
outside the current formalization.  Those claims retain the complete human
proofs and imported-source qualifications stated below.  ``New'' means new
relative to the pinned corpus; no unqualified worldwide priority claim is
made.  Conjectures and research problems are explicitly labelled and are not
used in the proof.
\end{statusbox}

\tableofcontents
\clearpage

\section{Main result and proof architecture}

\subsection{The theorem in one statement}

Throughout the report, $F:\R\to[0,1]$ denotes the bounded Fabius function:
it is $0$ on $(-\infty,0]$, $1$ on $[1,\infty)$, and on $[0,1]$ it is the
strictly increasing $C^\infty$ function satisfying
\begin{equation}\label{eq:F-basic}
 F(x)+F(1-x)=1,
 \qquad
 F'(x)=2F(2x)\quad(0\le x\le\tfrac12).
\end{equation}
Let $G:[0,1]\to[0,1]$ be its inverse, and define the totalized inverse
\begin{equation}\label{eq:totalized-inverse}
 \InvT(y):=G(\Clamp(y)),
 \qquad
 \Clamp(y):=\min(1,\max(0,y)).
\end{equation}
Thus $\InvT(y)=0$ for $y\le0$ and $\InvT(y)=1$ for $y\ge1$.

\begin{theorem}[Computability of the inverse Fabius function]\label{thm:main}
The inverse $G:[0,1]\to[0,1]$, in the subset-domain version of the
Grzegorczyk definition, and its totalization $\InvT:\R\to\R$ are computable
real functions.  More explicitly:
\begin{enumerate}
\item $G$ and $\InvT$ are sequentially computable.  From one algorithm that
      names a computable sequence $(y_i)$, one can uniformly compute dyadic
      approximations to $G(y_i)$ (or $\InvT(y_i)$) of any requested precision.
\item For $r\ge0$, put
      \begin{equation}\label{eq:Delta-intro}
       \Delta_r=\frac{1}{2^{r(r+1)/2}(r+1)^{r+1}}.
      \end{equation}
      Then, for all $u,v\in\R$,
      \begin{equation}\label{eq:explicit-dyadic-modulus-intro}
       \abs{u-v}<\Delta_r
       \quad\Longrightarrow\quad
       \abs{\InvT(u)-\InvT(v)}<2^{-r}.
      \end{equation}
\item For $n\ge1$, let
      \begin{equation}\label{eq:r-of-n}
       r(n):=\min\set{r\ge1:2^r>n},
      \end{equation}
      and define
      \begin{equation}\label{eq:d-of-n-intro}
       d(n):=2^{r(n)(r(n)+1)/2}\bigl(r(n)+1\bigr)^{r(n)+1},
       \qquad d(0):=1.
      \end{equation}
      The integer function $d$ is recursive (indeed primitive recursive after
      a standard bounded-search coding), and
      \begin{equation}\label{eq:wikipedia-modulus-intro}
       \abs{u-v}<\frac1{d(n)}
       \quad\Longrightarrow\quad
       \abs{\InvT(u)-\InvT(v)}<\frac1n
       \qquad(n\ge1).
      \end{equation}
\end{enumerate}
\end{theorem}

The theorem is stronger than a nonconstructive assertion that a continuous
inverse on a compact interval is uniformly continuous.  It gives a concrete
integer modulus, a finite algorithm on names, and a sharp structural identity
for the best possible modulus.

\subsection{Four proof layers}

The argument is organized so that every hidden effectiveness issue is exposed.

\begin{enumerate}
\item \textbf{Forward oracle.}  A centered finite convolution has an exact
Thue--Morse truncated-power formula.  At rational inputs it is rational, and
its error from $F$ is at most $2^{-N}$.
\item \textbf{Effective injectivity.}  Symmetric unimodality shows that an
interval of geometric length $h$ has probability mass at least $F(h)$;
strict radial unimodality shows that equality occurs only for a degenerate
interval or an endpoint interval.
A direct event gives the computable lower bound $F(2^{-r})\ge\Delta_r$.
\item \textbf{Uniform continuity.}  Effective injectivity of $F$ becomes an
effective modulus for $G$.  In fact the ordinary modulus of $G$ is exactly
$G$ itself.
\item \textbf{Name transformation.}  Tolerant bisection combines the forward
oracle and the inverse modulus.  A midpoint is used for a left/right update
only when the sign is certified; otherwise it is already a certified output.
\end{enumerate}

The fourth layer matters.  Naive bisection that waits until it decides whether
$F(c)<y$, $F(c)=y$, or $F(c)>y$ is not an algorithm on arbitrary computable
reals: equality of computable reals is undecidable in general.  The tolerance
branch is what removes that obstruction.

\section{Corpus audit, conventions, and nonduplication boundary}

\subsection{Arrival snapshot and live-corpus intake audit}

The delivered manuscript recorded commit \arrivalcommit\ \cite{ProveIt} as its
arrival snapshot and counted 96 LaTeX sources there.  On intake, the audit was
rerun against the live repository: the formal documentation audit scanned 582
modules and 7,976 public declarations, with no missing declaration comments or
module headers.  The present merged tree now scans 611 modules and 8,338
public declarations with the same zero-gap result.  A separate intake
inventory found 147 prior \texttt{.tex}
sources below \path{Analysis/FabiusFunction/docs}, excluding this newly filed
package.  These counts identify the corpus inspected; they are not a novelty
certificate.  All source paths were inventoried.  Primary and consolidated
mathematical reports were read directly; generated table fragments and frozen
variants were checked together with the parent documents that include them.

The inspected corpus spans the following overlapping families:
\begin{itemize}
\item the two source papers and the main Fabius/Rvachev exposition;
\item the Lean walkthrough, paper-coverage map, and mathematical glossary;
\item exact dyadic arithmetic, Thue--Morse and $q$-binomial formulae;
\item finite convolutions, infinite sinc products, Fourier decay, canonical
      products, and Poisson/Legendre representations;
\item endpoint asymptotics, lower-Lambert phases, Gamma--zeta periodic terms,
      Bell-polynomial inverse transfer, and all-orders inverse expansions;
\item inverse-dyadic sampling, $q$-Richardson acceleration, Lagrange and
      Legendre synthesis, total positivity, spectra, and operator viewpoints;
\item the final added spectral/$q$-Pochhammer and digital-geometry material.
\end{itemize}

The audit was thematic rather than a search for one spelling.  In particular,
it traced the bounded CDF, signed global extension, Rvachev function, totalized
inverse, finite spline evaluator, exact dyadic evaluator, inverse asymptotics,
and the repository's definitions of sequential computability and effective
uniform continuity.

\subsection{What the repository already proves}

The following inputs are established in the live development and are not
claimed as new here.

\begin{enumerate}
\item The bounded canonical $F$ exists, is unique, strictly increasing on
$[0,1]$, symmetric, $C^\infty$, and flat at both endpoints.
\item $F$ is the CDF of the random series in \cref{eq:random-series} below.
Its centered density is Rvachev's $\Up$ function, with Fourier transform an
infinite dyadic sinc product.
\item Every dyadic value of $F$ is an exact rational and is computed by an
executable recurrence.  The centered Thue--Morse splines converge uniformly
with error at most $2^{-N}$.
\item The forward bounded $F$ and the signed global extension are already
proved computable real functions.  In \path{FabiusComputability.lean} and
\path{FabiusComputableSpline.lean}, the public declarations include
\lean{abs_uniformCenteredPartialCDF_sub_fabiusReal_le},
\lean{abs_fabiusUniformSpline_sub_fabiusReal_le_all},
\lean{fabius_sequentiallyComputable}, and
\lean{fabius_isComputableRealFunction}.
\item \path{Convexity.lean} proves the strict rise and fall of the density via
\lean{strictMonoOn_deriv_fabiusReal_Icc} and
\lean{strictAntiOn_deriv_fabiusReal_Icc}; \path{Monotonicity.lean} proves
\lean{strictMonoOn_fabiusReal}.
\item \path{FabiusInverse.lean} constructs the clamped totalized inverse
\lean{fabiusInv}, proves the inverse identities
\lean{fabiusReal_fabiusInv} and \lean{fabiusInv_fabiusReal}, continuity via
\lean{continuous_fabiusInv}, and the interior differential and shape theory.
Exact dyadic inverse evaluation is also formalized, including
\lean{fabiusInv_fabiusAtInverseTwoPow}.
\item \path{FabiusInverseAsymptotic.lean} proves the leading equivalent
\lean{fabiusInv_isEquivalent_fabiusInverseAsymptoticMain}.  The companion
frontier reports develop explicit periodic corrections and all-orders
reversion, but those stronger formulas are not public Lean theorems.
\end{enumerate}

\subsection{The paper-level gap addressed here}

At the pinned arrival snapshot, no public theorem named or formulated as
sequential computability, effective uniform continuity, or
computable-real-function status for \lean{fabiusInv} occurred in the formal
corpus.  The missing bridge was not a formal triviality.  A
computable strictly increasing function can be inverted effectively on a
compact interval only after one supplies effective injectivity information;
ordinary strict monotonicity is qualitative.

This report supplies that information in two forms:
\begin{align}
 F(x+h)-F(x)&\ge F(h),\label{eq:gap-interval-preview}\\
 F(2^{-r})&\ge\Delta_r.\label{eq:gap-delta-preview}
\end{align}
The first is exact and geometric.  For $0<h\le1$, equality in the first bound
occurs exactly at $x=0$ or $x=1-h$; in the equivalent form
$F(b)-F(a)=F(b-a)$ the complete alternatives are $a=b$, $a=0$, or $b=1$.
The second bound is elementary and recursive.
Together with the existing forward evaluator they close both clauses of the
requested definition at the paper level.  Formalizing those clauses remains
the roadmap in \cref{sec:lean-status}.

The historical statement above remains about the pinned snapshot.  In the
current post-publication development, \path{InverseModulus.lean} formally
proves the first inequality, its complete equality locus, and its exact
transfer through the clamped inverse, including rigidity, sharpness, and
attainment.
\path{FabiusInverseEffectiveContinuity.lean} formally proves the second
\emph{inequality} by a stronger one-term recurrence estimate and packages a
simple primitive-recursive reciprocal modulus.
\path{FabiusInverseLogarithmicModulus.lean} now formalizes the report's smaller
logarithmic-$r(n)$ packaging, including its primitive-recursive selector and
denominators.  The report's probabilistic event proof and its
tolerant-bisection realizer remain on the human-proof side of the boundary
recorded in \cref{sec:lean-status}.

\subsection{Three functions that must not be conflated}

The documentation repeatedly warns about a notation trap.  The bounded CDF
$F$ is clamped outside $[0,1]$.  The signed global extension, often denoted
$\mathcal F$, agrees with $F$ on $[0,1]$ but oscillates outside it and obeys
global dilation/translation identities.  Rvachev's $\Up$ is an even compactly
supported bump on $[-1,1]$.  On its support,
\begin{equation}\label{eq:up-F-relation}
 \Up(t)=
 \begin{cases}
   F(t+1),&-1\le t\le0,\\
   F(1-t),&0\le t\le1.
 \end{cases}
\end{equation}
The inverse in this report is the inverse of the bounded transition
$F:[0,1]\to[0,1]$, with the clamped totalization \eqref{eq:totalized-inverse}.
It is not an inverse of the signed global extension.

\section{Computable-analysis definitions}

\subsection{Fast names and computable sequences}

Following standard computable analysis \cite{G1957,Weihrauch2000,PourElRichards1989}, a convenient representation of a real number $x$ is a fast dyadic Cauchy
name: an algorithm returns a dyadic rational $q_p$ such that
\begin{equation}\label{eq:fast-name}
 \abs{x-q_p}\le2^{-p}.
\end{equation}
A real sequence $(x_i)_{i\in\N}$ is computable when one algorithm returns
$q_{i,p}$ satisfying \eqref{eq:fast-name} uniformly in both $i$ and $p$.
The repository uses an equivalent signed-natural code for the dyadic
numerator, avoiding an opaque integer encoding.

\begin{definition}[Sequential computability]\label{def:sequential}
A function $f$ is \emph{sequentially computable} if every computable real
sequence $(x_i)$ in its domain is mapped to a computable sequence
$(f(x_i))$, uniformly in the sequence index and requested precision.
\end{definition}

The word ``uniformly'' is implicit in any useful implementation of this
definition: the output algorithm receives the input naming algorithm as a
subroutine.  The construction in \cref{sec:sequential} gives such a uniform
transformer, and therefore proves the weaker extensional statement as well.

\subsection{The reciprocal uniform-continuity convention}

\begin{definition}[Effective uniform continuity]\label{def:euc}
A function $f:\R\to\R$ is \emph{effectively uniformly continuous} if there
is a recursive $d:\N\to\N$ with $d(n)>0$ for $n>0$ such that
\begin{equation}\label{eq:euc-def}
 \abs{x-y}<\frac1{d(n)}
 \quad\Longrightarrow\quad
 \abs{f(x)-f(y)}<\frac1n
 \qquad(n>0).
\end{equation}
\end{definition}

This is the reciprocal convention used in the cited Wikipedia article
\cite{WikipediaCRF}.  The repository implements it in
\path{FabiusComputability.lean}.  It is equivalent to the more common binary
modulus convention
\begin{equation}\label{eq:binary-modulus-convention}
 \abs{x-y}<2^{-m(p)}
 \quad\Longrightarrow\quad
 \abs{f(x)-f(y)}<2^{-p},
\end{equation}
provided $m$ is recursive.  We give both forms because the binary version
reveals the endpoint conditioning much more clearly.

\begin{definition}[Computable real function]\label{def:computable-function}
In the Grzegorczyk formulation used here, a real function is computable if it
is sequentially computable and effectively uniformly continuous.
\end{definition}

For $G$ defined only on $[0,1]$, the same definitions are restricted to
sequences and pairs of points in that computable compact domain.  The totalized
$\InvT$ lets one use the literal type $\R\to\R$ without changing the
mathematics.

\section{Probability, Fourier structure, and the forward function}
\label{sec:forward}

\subsection{The random-series model}

The classical probability model may be written as follows\ \cite{Fabius1966,Arias05442,ProveIt}.  Let $U_1,U_2,\ldots$ be independent uniform random variables on $[0,1]$ and
put
\begin{equation}\label{eq:random-series}
 X:=\sum_{k=1}^{\infty}2^{-k}U_k.
\end{equation}
The series converges absolutely, $0\le X\le1$, and
\begin{equation}\label{eq:F-cdf}
 F(x)=\Prob[X\le x].
\end{equation}
Reflection $U_k\mapsto1-U_k$ sends $X$ to $1-X$ and proves the symmetry in
\eqref{eq:F-basic}.

The centered random variable
\begin{equation}\label{eq:centered-random}
 Y:=2X-1=\sum_{k=1}^{\infty}2^{-k}(2U_k-1)
\end{equation}
has density $\Up$.  With $\sinc z=\sin z/z$,
\begin{equation}\label{eq:sinc-product}
 \E\e^{\ii tY}
 =\prod_{k=1}^{\infty}\sinc\!\left(\frac{t}{2^k}\right).
\end{equation}
Thus the probability model, the Rvachev bump, and the infinite sinc product
are three representations of the same object.  The computability proof will
use the probability representation in physical space; \cref{sec:connections}
explains how the other representations can serve as alternative evaluators.

\subsection{Finite sums and their exact CDFs}

For $N\ge1$, define
\begin{equation}\label{eq:finite-random-sum}
 X_N:=\sum_{k=1}^{N}2^{-k}U_k,
 \qquad
 R_N:=X-X_N.
\end{equation}
Then
\begin{equation}\label{eq:tail-bound}
 0\le R_N\le2^{-N}.
\end{equation}
Let $F_N(x)=\Prob[X_N\le x]$.  The general inclusion--exclusion formula for a
sum of independent, nonidentically scaled uniforms becomes especially rigid
for dyadic scales.

\begin{proposition}[Finite Thue--Morse CDF]\label{prop:finite-cdf}
For $N\ge1$ and $x\in\R$,
\begin{equation}\label{eq:finite-cdf-formula}
 F_N(x)=
 \frac{2^{N(N+1)/2}}{N!}
 \sum_{j=0}^{2^N-1}(-1)^{s_2(j)}
 \pospart{x-\frac{j}{2^N}}^{N},
\end{equation}
where $s_2(j)$ is the sum of the binary digits of $j$.
\end{proposition}

\begin{proof}
For positive lengths $a_1,\ldots,a_N$, the CDF of
$\sum a_kU_k$ is
\[
 \frac{1}{N!\prod_{k=1}^N a_k}
 \sum_{A\subseteq\{1,\ldots,N\}}(-1)^{|A|}
 \pospart{x-\sum_{k\in A}a_k}^{N}.
\]
This follows by integrating the box indicator over the half-space
$\sum a_ku_k\le x$, or inductively by finite differences of truncated powers.
Here $a_k=2^{-k}$ and
$\prod a_k=2^{-N(N+1)/2}$.  The map
\[
 A\longmapsto j(A):=\sum_{k\in A}2^{N-k}
\]
is a bijection from subsets to $\{0,\ldots,2^N-1\}$, with
$\sum_{k\in A}2^{-k}=j(A)/2^N$ and
$|A|\equiv s_2(j(A))\pmod2$.  Substitution gives
\eqref{eq:finite-cdf-formula}.
\end{proof}

The alternating signs are exactly the Thue--Morse signs.  No limiting
argument is hidden in \eqref{eq:finite-cdf-formula}: at rational $x$ it is an
exact rational obtained by finitely many integer operations.

\subsection{Centered finite splines and a uniform certificate}

The deterministic midpoint of the omitted tail is $2^{-N-1}$.  Define
\begin{equation}\label{eq:centered-truncation}
 Q_N:=X_N+2^{-N-1},
 \qquad
 S_N(x):=\Prob[Q_N\le x]=F_N(x-2^{-N-1}).
\end{equation}
Then \cref{prop:finite-cdf} gives
\begin{equation}\label{eq:centered-spline-formula}
 S_N(x)=
 \frac{2^{N(N+1)/2}}{N!}
 \sum_{j=0}^{2^N-1}(-1)^{s_2(j)}
 \pospart{x-\frac{2j+1}{2^{N+1}}}^{N}.
\end{equation}
This is the centered uniform spline used by the repository's primitive-
recursive evaluator.

\begin{lemma}[Uniform density bound]\label{lem:density-bound}
The finite CDFs $F_N$ and $S_N$ are globally $2$-Lipschitz.
\end{lemma}

\begin{proof}
The density of $2^{-1}U_1$ is $2\mathbf1_{[0,1/2]}$.  Convolution with any
probability measure cannot increase the $L^\infty$ norm of a density.  Hence
the density of $X_N$ is bounded by $2$.  Translating by $2^{-N-1}$ does not
change this bound.  Integrating the density over an interval proves the
Lipschitz estimate.
\end{proof}

\begin{theorem}[Certified centered approximation]\label{thm:centered-error}
For every $N\ge1$ and every real $x$,
\begin{equation}\label{eq:centered-error}
 \abs{S_N(x)-F(x)}\le2^{-N}.
\end{equation}
\end{theorem}

\begin{proof}
Write
\[
 X=Q_N+E_N,
 \qquad
 E_N=R_N-2^{-N-1}.
\]
By \eqref{eq:tail-bound}, $|E_N|\le2^{-N-1}$.  Therefore
\[
 S_N(x-2^{-N-1})\le F(x)\le S_N(x+2^{-N-1}).
\]
By \cref{lem:density-bound}, moving the argument by $2^{-N-1}$ changes $S_N$
by at most $2\cdot2^{-N-1}=2^{-N}$.  This gives both sides of
\eqref{eq:centered-error}.
\end{proof}

\begin{corollary}[Self-contained computability of $F$]\label{cor:F-computable}
There is an explicit uniform algorithm which, given a rational $x$ and a
precision $p$, returns a rational $q$ with
$|q-F(x)|\le2^{-p}$.  One may take $q=S_N(x)$ with $N\ge p$.
Consequently $F$ is sequentially computable; together with its global
$2$-Lipschitz bound, it is a computable real function.
\end{corollary}

\begin{proof}
Formula \eqref{eq:centered-spline-formula} is a finite rational computation,
and \cref{thm:centered-error} supplies its error certificate.  At an arbitrary
computable input, first replace the input by a sufficiently accurate dyadic;
the $2$-Lipschitz bound controls this replacement.  The construction is
uniform in a sequence index.  Effective uniform continuity follows from
$|F(x)-F(y)|\le2|x-y|$; for example the reciprocal modulus $d(n)=2n$ works.
\end{proof}

\begin{remark}[Efficiency versus computability]
The direct sum has $2^N$ terms and is not intended as an efficient evaluator.
The repository's exact dyadic recurrence, highest-bit reduction, and
primitive-recursive natural-number implementation are substantially better.
For the theorem here, only termination and a proved error bound are required.
\end{remark}

\section{Effective positivity, strict increase, and the density shape}

\subsection{A closed positive lower bound at dyadic points}

The existence of an inverse requires strict increase; effective inversion
requires a quantitative version.  Both begin with an elementary positive
probability event.

\begin{lemma}[Dyadic endpoint box event]\label{lem:box-event}
For every $r\ge0$,
\begin{equation}\label{eq:Delta-def}
 F(2^{-r})\ge
 \Delta_r:=2^{-r(r+1)/2}(r+1)^{-(r+1)}.
\end{equation}
In particular, $F(x)>0$ for every $x>0$.
\end{lemma}

\begin{proof}
Put $m=r+1$.  For $1\le k\le m$, impose the independent restrictions
\begin{equation}\label{eq:box-restrictions}
 U_k\le a_k:=\frac{2^{k-r-1}}{m}.
\end{equation}
The largest $a_k$ is $a_m=1/m\le1$, so these are valid events.  On their
intersection,
\[
 2^{-k}U_k\le\frac{2^{-r-1}}m
 \quad(1\le k\le m).
\]
After summation the first $m$ terms contribute at most $2^{-r-1}$.  The
unrestricted tail contributes at most
\[
 \sum_{k=m+1}^{\infty}2^{-k}=2^{-m}=2^{-r-1}.
\]
Thus the event in \eqref{eq:box-restrictions} implies $X\le2^{-r}$.  Its
probability is
\begin{align*}
 \prod_{k=1}^{m}a_k
 &=m^{-m}2^{\sum_{k=1}^m(k-r-1)}\\
 &=m^{-m}2^{m(m+1)/2-m(r+1)}\\
 &=2^{-r(r+1)/2}(r+1)^{-(r+1)}=\Delta_r.
\end{align*}
This proves \eqref{eq:Delta-def}.  Given $x>0$, choose $r$ with
$2^{-r}\le x$ and use monotonicity of a CDF.
\end{proof}

The bound is deliberately elementary: it uses only independence, the support
of the tail, and integer arithmetic.  It is not numerically sharp, but
\cref{sec:binary-conditioning} shows that its logarithm has the correct two
largest orders for computability purposes.

The numerical conclusion of \cref{lem:box-event} is now Lean-formalized, but
by a different and stronger route.  Here is the complete recurrence argument.
Let $a_0=1$ and use the repository's exact inverse-dyadic identities
\[
 a_r=
 \frac{\displaystyle\sum_{0\le k<r}a_k/(r-k+1)!}{2^r-1}
 \quad(r\ge1),
 \qquad
 F(2^{-r})=2^{-\binom r2}a_r.
\]
Every $a_k$ is positive.  Retaining only the $k=0$ summand therefore gives
\[
 a_r\ge\frac{1}{(r+1)!(2^r-1)},
 \qquad
 F(2^{-r})\ge
 \frac{1}{2^{\binom r2}(r+1)!(2^r-1)}
 \quad(r\ge1).
\]
Since $2^r-1\le2^r$ and
$\binom r2+r=\binom{r+1}{2}$, taking the larger denominator yields the
first inequality below; the exceptional index $r=0$ is simply $F(1)=1$.
Thus for every $r\ge0$,
\begin{equation}\label{eq:factorial-endpoint-lower}
 F(2^{-r})\ge
 \frac{1}{2^{\binom{r+1}{2}}(r+1)!}
 \ge
 \frac{1}{2^{\binom{r+1}{2}}(r+1)^{r+1}}
 =\Delta_r.
\end{equation}
Indeed, $(r+1)!\le(r+1)^{r+1}$ proves the second inequality after
multiplication by the positive power of two and reversal under reciprocation.

For the executable arithmetic, write
\[
 d_{\mathrm{fac}}(r):=2^{\binom{r+1}{2}}(r+1)!,
 \qquad
 d_{\Delta}(r):=2^{\binom{r+1}{2}}(r+1)^{r+1}=\Delta_r^{-1}.
\]
Lean defines the first denominator by primitive recursion,
\[
 d_{\mathrm{fac}}(0)=1,
 \qquad
 d_{\mathrm{fac}}(r+1)
 =d_{\mathrm{fac}}(r)\,2^{r+1}(r+2).
\]
Induction, using
$\binom{r+2}{2}=\binom{r+1}{2}+r+1$, proves the displayed closed form.
The recursive clause uses only successor, multiplication, and exponentiation,
so $d_{\mathrm{fac}}$ is primitive recursive.  Likewise
$\binom{r+1}{2}=r(r+1)/2$, and primitive-recursive closure under successor,
multiplication, natural-number division, and exponentiation proves that
$d_{\Delta}$ is primitive recursive.  Finally
$d_{\mathrm{fac}}(r)\le d_{\Delta}(r)$ is exactly the factorial--power
inequality above, and positivity reverses this comparison between their
reciprocals.

The declarations
\path{Fabius.fabiusReal_inverse_two_pow_one_term_lower_bound},
\path{Fabius.inv_inverseFabiusFactorialDenominator_le_fabiusReal}, and
\path{Fabius.inv_inverseFabiusDeltaDenominator_le_fabiusReal} certify this
chain.  They do not formalize the event in
\eqref{eq:box-restrictions}, its containment, or its product probability;
that probabilistic proof remains report-level.  The argument in this paragraph
uses only the already formalized dyadic recurrence and therefore preserves the
distinction between the two proof routes.

\subsection{Strict monotonicity from the random series}

\begin{proposition}[Strict increase]\label{prop:strict-increase}
For $0\le a<b\le1$,
\begin{equation}\label{eq:strict-increase}
 F(a)<F(b).
\end{equation}
Hence $F:[0,1]\to[0,1]$ is a homeomorphism and $G=F^{-1}$ exists.
\end{proposition}

\begin{proof}
The density of $X_N$ is strictly positive on the interior of its support
$(0,1-2^{-N})$.  This follows inductively: a convolution of two densities
that are positive on the interiors of compact intervals is positive on the
interior of their Minkowski-sum interval, because the convolution integral
contains a nonempty open region on which both factors are positive.

Choose $N$ so large that $2^{-N}<b-a$.  Then the open interval
$(a,b-2^{-N})$ is nonempty and lies below the right endpoint of the support of
$X_N$.  Choose a nonempty open subinterval $J$ of it.  The preceding positivity
gives $\Prob[X_N\in J]>0$.  For $X_N\in J$ and every possible
$0\le R_N\le2^{-N}$, one has $a<X_N+R_N<b$.  Consequently
\[
 F(b)-F(a)=\Prob[a<X\le b]\ge\Prob[X_N\in J]>0.
\]
Continuity and the endpoint values $F(0)=0$, $F(1)=1$ then give the
homeomorphism statement.
\end{proof}

\subsection{Symmetric strict unimodality}

Let
\begin{equation}\label{eq:density-p}
 p(x):=F'(x),\qquad 0\le x\le1.
\end{equation}
The repository proves that $F$ is $C^\infty$, so $p$ is continuous.  The
functional and reflection equations completely determine its qualitative
shape.

\begin{proposition}[Shape of the Fabius density]\label{prop:density-shape}
The density $p$ satisfies
\begin{equation}\label{eq:density-symmetry}
 p(x)=p(1-x),
\end{equation}
is strictly increasing on $[0,1/2]$, and is strictly decreasing on
$[1/2,1]$.  Equivalently, $p(x)$ is a strictly decreasing function of
$|x-1/2|$.
\end{proposition}

\begin{proof}
On $[0,1/2]$, \eqref{eq:F-basic} gives
\[
 p(x)=2F(2x).
\]
By \cref{prop:strict-increase}, this is strictly increasing.  Differentiating
$F(x)+F(1-x)=1$ gives $p(x)=p(1-x)$, which reflects the first-half behavior
to the second half.
\end{proof}

This simple observation is the geometric heart of the inverse modulus.  It is
also a point where Rvachev's function enters directly: by
\eqref{eq:up-F-relation},
\begin{equation}\label{eq:p-up}
 p(x)=2\Up(2x-1).
\end{equation}
Thus \cref{prop:density-shape} is the standard single-peaked shape of the
compactly supported atomic bump.

\section{The least-mass interval and the exact inverse modulus}
\label{sec:exact-modulus}

\subsection{Intervals of a prescribed geometric length}

For $0\le h\le1$ and $0\le x\le1-h$, define the fixed-length CDF increment
\begin{equation}\label{eq:interval-mass}
 M_h(x):=F(x+h)-F(x).
\end{equation}
If $\mu$ is the Fabius probability law with CDF $F$, then the canonical
measure-theoretic interpretation is
\[
 M_h(x)=\mu\bigl((x,x+h]\bigr).
\]
The law is atomless---indeed it has the continuous density $p=F'$---so the
value is unchanged if either or both endpoints are included.  In particular,
\[
 \mu\bigl((x,x+h]\bigr)
 =\mu\bigl([x,x+h]\bigr)
 =\mu\bigl((x,x+h)\bigr)
 =\mu\bigl([x,x+h)\bigr).
\]
We therefore continue to call $M_h(x)$ an interval mass, while
\eqref{eq:interval-mass} fixes the canonical CDF convention.

\begin{theorem}[Strict interval-mass shape and exact minimizers]
\label{thm:least-mass}
For every $0\le h\le1$,
\begin{equation}\label{eq:least-mass}
 \min_{0\le x\le1-h}\bigl(F(x+h)-F(x)\bigr)=F(h).
\end{equation}
If $0<h\le1$ and $m=(1-h)/2$, then the globally extended increment $M_h$ is
strictly increasing on $[-h,m]$, symmetric about $m$, and strictly decreasing
on $[m,1]$; outside $[-h,1]$ it is identically zero.  On the admissible-start
interval $0\le x\le1-h$ one has the boundary-safe equality law
\begin{equation}\label{eq:least-mass-equality}
 M_h(x)=F(h)
 \quad\Longleftrightarrow\quad
 h=0\ \text{or}\ x=0\ \text{or}\ x=1-h.
\end{equation}
Equivalently, whenever $0\le a\le b\le1$,
\begin{equation}\label{eq:least-mass-equality-ab}
 F(b-a)=F(b)-F(a)
 \quad\Longleftrightarrow\quad
 a=b\ \text{or}\ a=0\ \text{or}\ b=1.
\end{equation}
Thus every genuinely interior interval $0<a<b<1$ has strictly more Fabius
mass than an endpoint interval of the same length.  At $h=1$ the two endpoint
formulas coincide at the single admissible start $x=0$.
\end{theorem}

\begin{proof}
The case $h=0$ gives $M_0\equiv0$, so suppose $0<h\le1$ and put
$m=(1-h)/2$.  Reflection gives
\begin{align*}
 M_h(1-h-x)
 &=F(1-x)-F(1-h-x)\\
 &=(1-F(x))-(1-F(x+h))=M_h(x),
\end{align*}
so $M_h$ is symmetric about $m$.

Consider $-h<x<m$.  Then $0<x+h<1$.  If $x\le0$, the constant left tail gives
$p(x)=0$, whereas strict positivity on the open support gives $p(x+h)>0$.
If $x>0$, compute
\begin{equation}\label{eq:distance-comparison}
 \left(x+h-\frac12\right)^2-
 \left(x-\frac12\right)^2
 =h(2x+h-1)<0.
\end{equation}
Thus $x+h$ is strictly closer to $1/2$ than $x$ is.  The strict radial
unimodality in \cref{prop:density-shape} again gives $p(x+h)>p(x)$.  Hence
\begin{equation}\label{eq:M-derivative}
 M_h'(x)=p(x+h)-p(x)>0
 \qquad(-h<x<m).
\end{equation}
The mean-value theorem makes $M_h$ strictly increasing on the closed interval
$[-h,m]$.  Reflection makes it strictly decreasing on $[m,1]$.  The clamped
tails of $F$ give $M_h=0$ to the left of $-h$ and to the right of $1$; this is
why the global half-line shape is only non-strict, whereas the two finite
branches above are maximally strict.

Finally,
\[
 M_h(0)=F(h)-F(0)=F(h)
 \quad\text{and}\quad
 M_h(1-h)=F(1)-F(1-h)=F(h).
\]
If $0<x<1-h$, compare $x$ strictly with $0$ on the left branch when $x\le m$,
and compare it strictly with $1-h$ on the right branch when $m<x$.  In either
case $M_h(x)>F(h)$.  This proves the minimum and the complete equality locus
\eqref{eq:least-mass-equality}.  Substituting $h=b-a$ and $x=a$ gives
\eqref{eq:least-mass-equality-ab}; its diagonal alternative $a=b$ is exactly
the degenerate case $h=0$.
\end{proof}

\begin{corollary}[Global increment shape and constrained superadditivity]
\label{cor:global-increment-shape}
For every $h\ge0$, the globally extended increment $M_h$ is nondecreasing on
the half-line
\[
 x\le\frac{1-h}{2}
\]
and nonincreasing on the complementary half-line.  Consequently, whenever
$a,b\ge0$ and $a+b\le1$,
\begin{equation}\label{eq:fabius-superadditive}
 F(a)+F(b)\le F(a+b).
\end{equation}
Its complete equality locus is
\begin{equation}\label{eq:fabius-superadditive-equality}
 F(a)+F(b)=F(a+b)
 \quad\Longleftrightarrow\quad
 a=0\ \text{or}\ b=0\ \text{or}\ a+b=1.
\end{equation}
\end{corollary}

\begin{proof}
The derivative argument above never used $x\ge0$: from
$x\le(1-h)/2$ and $h\ge0$, either $x+h\le1/2$ and the monotonicity of $p$ on
the left half gives $p(x)\le p(x+h)$, or $x+h>1/2$ and reflection replaces
$p(x+h)$ by $p(1-x-h)$, where
\[
 x\le1-x-h\le\frac12.
\]
Thus $M_h'\ge0$ on the entire left half-line.  Reflection
$M_h(1-h-x)=M_h(x)$ gives the right-half statement.  For
\eqref{eq:fabius-superadditive}, apply \cref{thm:least-mass} to the interval
from $a$ to $a+b$: its length is $b$, so
\[
 F(b)\le F(a+b)-F(a).
\]
 Rearranging proves the inequality.  Equality is exactly the equality case of
 \cref{thm:least-mass} for the interval $[a,a+b]$: its three alternatives are
 $a=a+b$, $a=0$, and $a+b=1$, and the first is equivalent to $b=0$.
\end{proof}

\begin{remark}[A general principle]\label{rem:general-principle}
The lower bound used only that a probability density on $[0,1]$ is symmetric
and unimodal.  Hence the minimum-value part of \cref{thm:least-mass} holds for
every continuous symmetric unimodal distribution.  The endpoint-only equality
locus uses the genuinely stronger hypothesis that the density decreases
strictly with distance from the midpoint; plateaus can create additional
minimizers.  The Fabius system is special not because the geometric lemma is
exotic, but because $F(h)$ has unusually rich exact arithmetic and endpoint
asymptotics.
\end{remark}

\subsection{The inverse's modulus is the inverse itself}

For a bounded function $H$ define its ordinary modulus of continuity by
\begin{equation}\label{eq:ordinary-modulus}
 \omega_H(\eta):=
 \sup\set{\abs{H(u)-H(v)}:\abs{u-v}\le\eta}.
\end{equation}
For $G$ the variables are restricted to $[0,1]$; for $\InvT$ they range over
$\R$.

\begin{theorem}[Exact modulus of the inverse]\label{thm:exact-inverse-modulus}
For every $\eta\ge0$, with the inputs in the definition of $\omega_G$
restricted to $[0,1]$,
\begin{equation}\label{eq:exact-inverse-modulus}
 \omega_G(\eta)=G(\min(\eta,1)).
\end{equation}
In particular this is $G(\eta)$ for $0\le\eta\le1$ and is $1$ thereafter.
For the totalized inverse and every $\eta\ge0$,
\begin{equation}\label{eq:exact-total-modulus}
 \omega_{\InvT}(\eta)=G(\min(\eta,1)).
\end{equation}
More strongly, for all real $u,v$,
\begin{equation}\label{eq:absolute-inverse-gap-bound}
 \abs{\InvT(u)-\InvT(v)}
 \le \InvT(\abs{u-v})
 =G\bigl(\min(\abs{u-v},1)\bigr).
\end{equation}
The associated ordered-gap and subadditivity laws are
\begin{equation}\label{eq:inverse-gap-bound}
 \InvT(v)-\InvT(u)\le \InvT(v-u)
 \qquad(u,v\in\R)
\end{equation}
and
\begin{equation}\label{eq:inverse-subadditive}
 \InvT(x+y)\le\InvT(x)+\InvT(y)
 \qquad(x,y\in\R).
\end{equation}
On the unit interval the ordered-gap equality cases are exactly
\begin{equation}\label{eq:inverse-gap-equality}
 G(v)-G(u)=G(v-u)
 \quad\Longleftrightarrow\quad
 u=v\ \text{or}\ u=0\ \text{or}\ v=1
 \qquad(0\le u\le v\le1).
\end{equation}
Equivalently, without ordering the inputs,
\begin{equation}\label{eq:absolute-inverse-gap-equality}
 \abs{G(u)-G(v)}=G(\abs{u-v})
 \quad\Longleftrightarrow\quad
 u=v\ \text{or}\ u\in\{0,1\}\ \text{or}\ v\in\{0,1\}
 \qquad(u,v\in[0,1]).
\end{equation}
These two equality classifications are intentionally unit-domain statements:
the constant tails of $\InvT$ create additional configurations on $\R$.
\end{theorem}

\begin{proof}
First let $c,d\in[0,1]$.  Interchanging them if necessary, assume $c\le d$,
and put
\[
 a=G(c),\qquad b=G(d).
\]
Monotonicity of $G$ gives $a\le b$.  The inverse identities and
\cref{thm:least-mass} give
\begin{equation}\label{eq:mass-gap-proof}
 F(b-a)\le F(b)-F(a)=d-c.
\end{equation}
Both $b-a$ and $d-c$ lie in $[0,1]$.  Applying the increasing inverse and
using $G(F(b-a))=b-a$ yields
\[
 b-a\le G(d-c).
\]
Removing the temporary ordering assumption gives
\begin{equation}\label{eq:unit-absolute-gap}
 \abs{G(c)-G(d)}\le G(\abs{c-d})
 \qquad(c,d\in[0,1]).
\end{equation}
The same argument records equality.  Under $c\le d$, equality in the inverse
gap is equivalent, after applying the strictly increasing $F$, to equality in
\eqref{eq:mass-gap-proof}.  By \eqref{eq:least-mass-equality-ab}, this says
\[
 a=b,\qquad a=0,\qquad\text{or}\qquad b=1.
\]
Because $a=G(c)$ and $b=G(d)$ and $F,G$ are inverse order isomorphisms, these
alternatives are respectively $c=d$, $c=0$, and $d=1$.  This proves
\eqref{eq:inverse-gap-equality}.  Applying it after ordering $c,d$ proves
\eqref{eq:absolute-inverse-gap-equality}; reversing the order exchanges the
left and right endpoint alternatives.

Now put $c=\Clamp(u)$ and $d=\Clamp(v)$.  Projection onto $[0,1]$ is
nonexpansive and has diameter one, so
\[
 \abs{c-d}\le\min(\abs{u-v},1).
\]
Combining this with \eqref{eq:unit-absolute-gap} proves
\eqref{eq:absolute-inverse-gap-bound}.  If $\abs{u-v}\le\eta$, monotonicity
then gives the upper bound $G(\min(\eta,1))$ for both moduli.  Let
$\theta=\min(\eta,1)$.  The pair $(0,\theta)$ belongs to $[0,1]^2$, has input
separation $\theta\le\eta$, and has output separation
\[
 \abs{G(\theta)-G(0)}=G(\theta).
\]
Thus the displayed upper bound is attained, not merely approached, proving
both exact supremum identities.

For \eqref{eq:inverse-gap-bound}, if $u\le v$, remove the absolute value from
\eqref{eq:absolute-inverse-gap-bound}; then $v-u=\abs{v-u}$.  If $v<u$, the
left side is nonpositive by monotonicity, whereas $\InvT(v-u)\ge0$.
Finally apply \eqref{eq:inverse-gap-bound} with $u=x$ and $v=x+y$:
\[
 \InvT(x+y)-\InvT(x)\le\InvT(y),
\]
which rearranges to \eqref{eq:inverse-subadditive}.
\end{proof}

\begin{corollary}[Effective-injectivity form]\label{cor:effective-injectivity}
For $0\le h\le1$ and $u,v\in\R$,
\begin{equation}\label{eq:effective-injectivity}
 \abs{u-v}<F(h)
 \quad\Longrightarrow\quad
 \abs{\InvT(u)-\InvT(v)}<h.
\end{equation}
The closed-threshold and contrapositive forms are
\begin{align}
 \abs{u-v}\le F(h)
 &\Longrightarrow \abs{\InvT(u)-\InvT(v)}\le h,
 \label{eq:closed-effective-injectivity}\\
 h\le\abs{\InvT(u)-\InvT(v)}
 &\Longrightarrow F(h)\le\abs{u-v}.
 \label{eq:contrapositive-effective-injectivity}
\end{align}
The same statements hold for $G$ on $[0,1]$.
\end{corollary}

\begin{proof}
The strict hypothesis gives
$\abs{u-v}<F(h)\le1$, so the saturation in the exact modulus disappears.
Strict monotonicity of $G$ on $[0,1]$ therefore gives
\[
 \abs{\InvT(u)-\InvT(v)}
 \le G(\abs{u-v})<G(F(h))=h.
\]
The same chain with non-strict inequalities proves
\eqref{eq:closed-effective-injectivity}; its contrapositive is
\eqref{eq:contrapositive-effective-injectivity}.
\end{proof}

\begin{theorem}[Exact strict threshold]\label{thm:exact-strict-threshold}
For $0\le h\le1$ and every real $\delta$,
\begin{equation}\label{eq:exact-strict-threshold}
 \left[
   \forall u,v\in\R,
   \ \abs{u-v}<\delta
   \Longrightarrow
   \abs{\InvT(u)-\InvT(v)}<h
 \right]
 \quad\Longleftrightarrow\quad
 \delta\le F(h).
\end{equation}
Thus $F(h)$ is exactly the largest strict universal input threshold for output
error below $h$.
\end{theorem}

\begin{proof}
If $\delta\le F(h)$, then
$\abs{u-v}<\delta$ implies $\abs{u-v}<F(h)$, and
\cref{cor:effective-injectivity} applies.  Conversely, suppose the universal
implication holds but $F(h)<\delta$.  Take
\[
 u=0,\qquad v=F(h).
\]
Their input separation is $F(h)<\delta$, whereas the inverse identities give
\[
 \abs{\InvT(0)-\InvT(F(h))}=\abs{0-h}=h,
\]
contradicting the asserted strict output bound.  Notice that no sign
hypothesis on $\delta$ is needed: for $\delta\le0$ the left implication is
vacuous and the right inequality follows from $F(h)\ge0$.
\end{proof}

\section{Effective uniform continuity}

\subsection{A closed primitive-recursive modulus}

Combining the exact inverse modulus with the box event gives the promised
fully explicit implication.

\begin{theorem}[Closed dyadic modulus]\label{thm:closed-dyadic-modulus}
For every $r\ge0$ and all $u,v\in\R$,
\begin{equation}\label{eq:closed-dyadic-modulus}
 \abs{u-v}<\Delta_r
 \quad\Longrightarrow\quad
 \abs{\InvT(u)-\InvT(v)}<2^{-r},
\end{equation}
where $\Delta_r$ is given by \eqref{eq:Delta-def}.
\end{theorem}

\begin{proof}
By \cref{lem:box-event}, $\Delta_r\le F(2^{-r})$.  Apply
\cref{cor:effective-injectivity} with $h=2^{-r}$.
\end{proof}

The exact statement \eqref{eq:closed-dyadic-modulus} is also formalized as
\path{Fabius.abs_fabiusInv_sub_lt_inverse_two_pow_of_lt_deltaDenominator}.
The stronger factorial-denominator version is
\path{Fabius.abs_fabiusInv_sub_lt_inverse_two_pow_of_lt_factorialDenominator}.
The formal API also closes both endpoints:
\[
 \abs{u-v}\le\Delta_r
 \quad\Longrightarrow\quad
 \abs{\InvT(u)-\InvT(v)}\le2^{-r},
\]
with declaration
\path{Fabius.abs_fabiusInv_sub_le_inverse_two_pow_of_le_deltaDenominator};
\path{Fabius.abs_fabiusInv_sub_le_inverse_two_pow_of_le_factorialDenominator}
is its stronger factorial-denominator companion.  Thus the strict and closed
input/output boundaries are separately available rather than recovered by an
unstated limiting argument.  Their common proof is direct.  Let $d$ be either
$d_{\mathrm{fac}}(r)$ or $d_{\Delta}(r)$.  The recurrence argument above gives
$d^{-1}\le F(2^{-r})$.  Hence
\[
 \abs{u-v}<d^{-1}
 \quad\Longrightarrow\quad
 \abs{u-v}<F(2^{-r})
 \quad\Longrightarrow\quad
 \abs{\InvT(u)-\InvT(v)}<2^{-r}
\]
by strict effective injectivity.  Replacing both strict inequalities by weak
ones and using its closed-threshold companion proves the two closed versions.
This proves all four formal dyadic implications without a limiting argument;
only the choice of denominator distinguishes the factorial and Delta forms.

\begin{theorem}[Reciprocal modulus required by the definition]\label{thm:reciprocal-modulus}
For $n\ge1$, let $r(n)$ be the least positive integer with $2^{r(n)}>n$ and
put
\begin{equation}\label{eq:d-explicit}
 d(n)=2^{r(n)(r(n)+1)/2}\bigl(r(n)+1\bigr)^{r(n)+1},
 \qquad d(0)=1.
\end{equation}
Then $d$ is recursive and satisfies \eqref{eq:euc-def} for $\InvT$.
Thus $\InvT$ and the restricted inverse $G$ are effectively uniformly
continuous.
\end{theorem}

\begin{proof}
The function $r(n)$ can be found by bounded integer search, or equivalently by
computing the binary length of $n$; hence it is recursive.  Exponentiation,
multiplication, and composition preserve recursiveness, so $d$ is recursive.
By construction,
\[
 \frac1{d(n)}=\Delta_{r(n)}
 \quad\text{and}\quad
 2^{-r(n)}<\frac1n.
\]
If $|u-v|<1/d(n)$, \cref{thm:closed-dyadic-modulus} gives
$|\InvT(u)-\InvT(v)|<2^{-r(n)}<1/n$.
\end{proof}

\begin{remark}[An even simpler, much larger modulus]
Taking $r=n$ gives the one-line primitive-recursive choice
\[
 d_0(n)=2^{n(n+1)/2}(n+1)^{n+1}\qquad(n\ge1),
\]
because $2^{-n}<1/n$.  The logarithmic choice $r(n)$ in
\cref{thm:reciprocal-modulus} is far smaller and has the correct leading
complexity scale.

The Lean development uses the stronger one-line witness
\[
 d_{\mathrm{fac}}(n)=2^{\binom{n+1}{2}}(n+1)!,
\]
which is no larger than $d_0(n)$.  The theorem
\path{Fabius.fabiusInv_effectivelyUniformContinuous} proves
\texttt{EffectivelyUniformContinuous} for the totalized inverse as follows.
The recurrence above makes $d_{\mathrm{fac}}$ primitive recursive and positive.
For $n\ge1$, the strict factorial implication gives
\[
 \abs{u-v}<\frac1{d_{\mathrm{fac}}(n)}
 \quad\Longrightarrow\quad
 \abs{\InvT(u)-\InvT(v)}<2^{-n}.
\]
The elementary inequality $n<2^n$ between positive integers reverses under
reciprocation, so $2^{-n}<1/n$.  Chaining these inequalities is precisely the
reciprocal convention in \cref{def:euc}; $d_{\mathrm{fac}}(0)=1$ supplies the
harmless zeroth value.  Thus effective uniform continuity itself is
formalized.  The downstream
\path{FabiusInverseLogarithmicModulus.lean} now also formalizes the smaller
choice $r(n)$ in \cref{thm:reciprocal-modulus}.  Its declaration
\path{Fabius.inverseFabiusLogarithmicOrder_isLeast} identifies the least order,
\path{Fabius.inverseFabiusLogarithmicOrder_primrec} proves primitive
recursiveness, and
\path{Fabius.fabiusInv_effectivelyUniformContinuous_logarithmicDelta} packages
the report-exact denominator $d(n)$.  The same module supplies the stronger
logarithmic factorial witness.
\end{remark}

\subsection{The exact-dyadic repository modulus}

Put $h_n:=2^{-r(n)}$.  The corpus proves that every $F(h_n)$ is an exactly
computable positive rational.  Therefore one can replace the elementary lower
bound $\Delta_{r(n)}$ by the exact endpoint mass while retaining an exact
dyadic output proxy.

\begin{proposition}[Repository-native exact-dyadic modulus]\label{prop:exact-dyadic-modulus}
Let
\begin{equation}\label{eq:d-star}
 d_*(n):=
 \ceil{\frac{1}{F(h_n)}}
 \qquad(n\ge1),
 \qquad h_n=2^{-r(n)},
 \quad d_*(0):=1.
\end{equation}
Then $d_*$ is recursive and
\begin{equation}\label{eq:d-star-strong}
 \abs{u-v}<\frac1{d_*(n)}
 \quad\Longrightarrow\quad
 \abs{\InvT(u)-\InvT(v)}<h_n<\frac1n.
\end{equation}
Moreover, $d_*(n)$ is the least positive integer $d$ satisfying
\[
 \frac1d\le F(h_n).
\]
Equivalently, it is the least reciprocal denominator certified by the exact
inverse-modulus law for the fixed dyadic output target $h_n$.  It is not
claimed to be the least integer modulus for the weaker target $1/n$.
\end{proposition}

\begin{proof}
The exact dyadic evaluator computes the positive rational $F(h_n)$, and
rational reciprocal and ceiling are computable integer operations.  By the
defining property of the ceiling,
\[
 \frac1{d_*(n)}\le F(h_n).
\]
The strict input hypothesis and \cref{cor:effective-injectivity} therefore
give output separation strictly below $h_n$, while the definition of $r(n)$
gives $h_n<1/n$.

For the qualified minimality statement, let $1\le d<d_*(n)$.  Then
\[
 d<\frac1{F(h_n)},
 \qquad\text{hence}\qquad
 \frac1d>F(h_n).
\]
The endpoint pair $u=0$, $v=F(h_n)$ has input separation strictly below
$1/d$ and output separation exactly $h_n$.  Thus $d$ cannot certify the
stricter conclusion $\abs{\InvT(u)-\InvT(v)}<h_n$.
\end{proof}

\begin{remark}[The genuine integer optimum for the $1/n$ target]
The exact inverse-modulus law shows that a positive integer $d$ has the
universal property
\[
 \abs{u-v}<\frac1d
 \quad\Longrightarrow\quad
 \abs{\InvT(u)-\InvT(v)}<\frac1n
\]
if and only if
\[
 \frac1d\le F(1/n).
\]
Consequently the mathematically least such denominator is
\[
 d_{\mathrm{opt}}(n)=\ceil{\frac1{F(1/n)}}.
\]
The present report does not prove that $n\mapsto d_{\mathrm{opt}}(n)$ is
recursive: computability of the real number $F(1/n)$ alone does not make the
discontinuous ceiling operation effective at an unknown integer boundary.
The dyadic quantity $d_*$ avoids that issue because $F(h_n)$ is an exact
rational.

The distinction is already visible at $n=1$.  Here $r(1)=1$,
$F(1/2)=1/2$, and hence $d_*(1)=2$, whereas $d=1$ already satisfies the
required implication with output tolerance $1$.
\end{remark}

The exact-ceiling construction $d_*$, including its recursion and qualified
denominator minimality, has not yet been translated into Lean.  The current
logarithmic effective-continuity layer deliberately avoids the ceiling by
using the explicit Delta or stronger factorial denominator at the formalized
least order $r(n)$.

\subsection{Binary conditioning is quadratic}
\label{sec:binary-conditioning}

The exact modulus permits a sharp, representation-independent measurement of
how many input bits are intrinsically needed near the flat endpoint.

\begin{definition}[Optimal binary modulus index]\label{def:mu}
For $p\ge1$, define
\begin{equation}\label{eq:mu-def}
 \mu(p):=\min\set{m\in\N:2^{-m}\le F(2^{-p})}
 =\ceil{\log_2\frac1{F(2^{-p})}}.
\end{equation}
Then
\begin{equation}\label{eq:mu-property}
 \abs{u-v}<2^{-\mu(p)}
 \quad\Longrightarrow\quad
 \abs{\InvT(u)-\InvT(v)}<2^{-p},
\end{equation}
and no smaller $m$ has this universal property.
\end{definition}

The lower event bound already gives an upper bound for $\mu(p)$.  A matching
lower bound follows from endpoint flatness.

\begin{lemma}[Elementary flatness upper bound]\label{lem:flatness-upper}
For every integer $q\ge1$ and $0\le x\le2^{-q}$,
\begin{equation}\label{eq:flatness-upper}
 F(x)\le\frac{2^{q(q+1)/2}}{q!}\,x^q.
\end{equation}
In particular,
\begin{equation}\label{eq:dyadic-flatness-upper}
 F(2^{-p})\le\frac{2^{-p(p-1)/2}}{p!}.
\end{equation}
\end{lemma}

\begin{proof}
Repeated differentiation of $F'(x)=2F(2x)$ gives
\begin{equation}\label{eq:iterated-derivative}
 F^{(q)}(x)=2^{q(q+1)/2}F(2^qx)
 \qquad(0\le x\le2^{-q}).
\end{equation}
All derivatives of $F$ vanish at $0$, and $0\le F\le1$.  Taylor's formula in
integral form therefore gives
\begin{align*}
 F(x)
 &=\frac1{(q-1)!}\int_0^x(x-t)^{q-1}F^{(q)}(t)\dd t\\
 &\le\frac{2^{q(q+1)/2}}{(q-1)!}
      \int_0^x(x-t)^{q-1}\dd t
 =\frac{2^{q(q+1)/2}}{q!}x^q.
\end{align*}
Set $q=p$ and $x=2^{-p}$ to obtain
\eqref{eq:dyadic-flatness-upper}.
\end{proof}

\begin{theorem}[Asymptotically sharp input-bit requirement]\label{thm:mu-asymptotic}
For every $p\ge1$,
\begin{align}
 \ceil{\frac{p(p-1)}2+\log_2(p!)}
 \;&\le\;\mu(p)\label{eq:mu-lower}\\
 &\le\;
 \ceil{\frac{p(p+1)}2+(p+1)\log_2(p+1)}.
 \label{eq:mu-upper}
\end{align}
Consequently
\begin{equation}\label{eq:mu-asymptotic}
 \boxed{\displaystyle
 \mu(p)=\frac{p^2}{2}+p\log_2p+\bigO(p).}
\end{equation}
\end{theorem}

\begin{proof}
The flatness bound \eqref{eq:dyadic-flatness-upper} yields
\[
 \log_2\frac1{F(2^{-p})}
 \ge\frac{p(p-1)}2+\log_2(p!),
\]
while \cref{lem:box-event} yields
\[
 \log_2\frac1{F(2^{-p})}
 \le\frac{p(p+1)}2+(p+1)\log_2(p+1).
\]
Take ceilings and use \eqref{eq:mu-def}.  Stirling's formula
$\log_2(p!)=p\log_2p-(\log_2\e)p+\bigO(\log p)$ makes both bounds equal to
$p^2/2+p\log_2p+\bigO(p)$.
\end{proof}

\begin{remark}[Intrinsic endpoint cost]
The quadratic precision loss is not an artifact of the spline evaluator or of
bisection.  The two admissible inputs $0$ and $F(2^{-p})$ have outputs separated
by $2^{-p}$.  Any uniform oracle algorithm that promises $p$ output bits must,
in the worst endpoint case, obtain enough input information to distinguish
those values.  The exact modulus quantifies that information barrier.
\end{remark}

\begin{corollary}[Growth of the reciprocal modulus]\label{cor:d-growth}
For $n\to\infty$, an asymptotically scale-optimal reciprocal modulus has
\begin{equation}\label{eq:d-growth}
 \log_2 d(n)
 =\frac12(\log_2 n)^2
  +(\log_2n)\log_2\log_2n
  +\bigO(\log n).
\end{equation}
Equivalently,
\begin{equation}\label{eq:d-quasipoly}
 d(n)=n^{\frac12\log_2n+\log_2\log_2n+\bigO(1)}.
\end{equation}
\end{corollary}

\begin{proof}
Use $p=r(n)=\log_2n+\bigO(1)$ in
\cref{thm:mu-asymptotic}.  The elementary modulus of
\cref{thm:reciprocal-modulus} obeys the same two leading orders directly.
\end{proof}

\section{Sequential computability by tolerant bisection}
\label{sec:sequential}

\subsection{Why ordinary bisection is insufficient}

Suppose $y$ is given only by a computable name.  At a dyadic midpoint $c$,
one can approximate both $F(c)$ and $y$, but in general no algorithm can
decide whether the two represented reals are exactly equal.  An implementation
that loops until it proves either $F(c)<y$ or $F(c)>y$ can therefore fail to
terminate when $y=F(c)$.

The inverse modulus provides the missing third branch.  If the approximations
are too close to certify a sign, then $F(c)$ is close to $y$; effective
injectivity says that $c$ is close to $G(y)$, so the algorithm may return it.

\subsection{The algorithm on a computable sequence}

Let $(y_i)$ be a computable sequence in $[0,1]$.  Fix a naming algorithm
$q(i,s)\in\Q$ with
\begin{equation}\label{eq:input-name}
 \abs{q(i,s)-y_i}\le2^{-s}.
\end{equation}
We describe an algorithm which, from $(i,p)$, returns a dyadic $z(i,p)$ with
\begin{equation}\label{eq:output-name-goal}
 \abs{z(i,p)-G(y_i)}<2^{-p}.
\end{equation}

\begin{algorithm}[Tolerant certified bisection]\label{alg:tolerant-bisection}
Given an index $i$ and output precision $p$:
\begin{enumerate}
\item Set
\[
 r=p+2,
 \qquad h=2^{-r},
 \qquad \delta=\Delta_r.
\]
Choose integers $s,N$ effectively so that
\begin{equation}\label{eq:oracle-precisions}
 2^{-s}\le\delta/8,
 \qquad 2^{-N}\le\delta/8.
\end{equation}
Compute the rational $\widetilde y=q(i,s)$.
\item Initialize the certified bracket $[a,b]=[0,1]$.
\item Repeat at most $p+2$ times:
  \begin{enumerate}
  \item Put $c=(a+b)/2$ and evaluate the exact rational $S_N(c)$ by
        \eqref{eq:centered-spline-formula}.
  \item Form the rational approximate difference
        \begin{equation}\label{eq:D-tilde}
         \widetilde D=S_N(c)-\widetilde y.
        \end{equation}
  \item If $\widetilde D>\delta/2$, replace $b$ by $c$.
  \item If $\widetilde D<-\delta/2$, replace $a$ by $c$.
  \item Otherwise return $c$.
  \end{enumerate}
\item If no midpoint was returned, return $(a+b)/2$ after the final update.
\end{enumerate}
\end{algorithm}

Every comparison in \cref{alg:tolerant-bisection} is between exact rationals.
The order $N$ can be extremely large, but it is obtained by a finite integer
calculation and the spline sum is finite.  Hence the procedure is an algorithm
in the computability-theoretic sense.

\subsection{Correctness of every branch}

\begin{lemma}[Difference certificate]\label{lem:difference-certificate}
At every midpoint $c$ in \cref{alg:tolerant-bisection},
\begin{equation}\label{eq:difference-certificate}
 \abs{\widetilde D-(F(c)-y_i)}\le\delta/4.
\end{equation}
\end{lemma}

\begin{proof}
By \cref{thm:centered-error} and \eqref{eq:oracle-precisions},
$|S_N(c)-F(c)|\le\delta/8$.  By \eqref{eq:input-name} and the same choice of
$s$, $|\widetilde y-y_i|\le\delta/8$.  The triangle inequality gives
\eqref{eq:difference-certificate}.
\end{proof}

\begin{lemma}[Safe left/right updates]\label{lem:safe-updates}
If $\widetilde D>\delta/2$, then $G(y_i)<c$.  If
$\widetilde D<-\delta/2$, then $c<G(y_i)$.
\end{lemma}

\begin{proof}
In the first case, \cref{lem:difference-certificate} gives
\[
 F(c)-y_i\ge\widetilde D-\delta/4>\delta/4>0.
\]
Since $F$ is strictly increasing, $c>G(y_i)$.  The second case is symmetric.
Thus every update preserves the invariant
\begin{equation}\label{eq:bracket-invariant}
 G(y_i)\in[a,b].
\end{equation}
\end{proof}

\begin{lemma}[The inconclusive branch is a success branch]\label{lem:near-branch}
If $|\widetilde D|\le\delta/2$, then
\begin{equation}\label{eq:near-output}
 \abs{c-G(y_i)}<h=2^{-r}<2^{-p}.
\end{equation}
\end{lemma}

\begin{proof}
By \cref{lem:difference-certificate},
\[
 |F(c)-y_i|\le|\widetilde D|+\delta/4
 \le3\delta/4<\delta.
\]
The choice $\delta=\Delta_r$ and \cref{thm:closed-dyadic-modulus}, applied to
the input pair $F(c),y_i$, give
\[
 |G(F(c))-G(y_i)|<2^{-r}.
\]
Because $c\in[0,1]$, $G(F(c))=c$.  Finally, $r=p+2$.
\end{proof}

\begin{theorem}[Correctness and uniformity of tolerant bisection]
\label{thm:bisection-correct}
Algorithm \ref{alg:tolerant-bisection} terminates and returns a dyadic
$z(i,p)$ satisfying \eqref{eq:output-name-goal}.  The map $(i,p)\mapsto z(i,p)$
is computable whenever the name $(i,s)\mapsto q(i,s)$ is computable.
\end{theorem}

\begin{proof}
The loop has a fixed finite bound.  If it returns through the inconclusive
branch, \cref{lem:near-branch} proves the error estimate.  Otherwise
\cref{lem:safe-updates} preserves the bracket invariant through $p+2$
bisections.  Its final width is $2^{-(p+2)}$, so the final midpoint is at
distance at most $2^{-(p+3)}<2^{-p}$ from $G(y_i)$.

All quantities $r,h,\delta,s,N$ are computable from $p$.  Midpoints are
dyadic; $S_N(c)$ is an exact finite rational expression; and every branch is
a rational comparison.  The only access to the sequence is the single call to
its naming algorithm.  Therefore the construction is uniform in $i$ and $p$.
\end{proof}

\begin{corollary}[Sequential computability on the unit interval]
\label{cor:G-sequential}
The inverse $G:[0,1]\to[0,1]$ is sequentially computable.
\end{corollary}

\begin{proof}
Given any computable sequence $(y_i)$ in $[0,1]$, the output of
\cref{alg:tolerant-bisection} is a computable fast dyadic name for the sequence
$(G(y_i))$.
\end{proof}

\subsection{Clamping and the totalized inverse}

\begin{lemma}[Computability of clamping]\label{lem:clamp-computable}
If $(y_i)$ is a computable real sequence, then
$(\Clamp(y_i))$ is a computable sequence.
\end{lemma}

\begin{proof}
Clamping is $1$-Lipschitz.  If $q(i,p)$ approximates $y_i$ within $2^{-p}$,
then the exactly rational quantity $\Clamp(q(i,p))$ approximates
$\Clamp(y_i)$ within the same error.  The rational min/max operations are
decidable and computable.
\end{proof}

\begin{corollary}[Sequential computability of the totalized inverse]
\label{cor:GT-sequential}
The totalized inverse $\InvT:\R\to\R$ is sequentially computable.
\end{corollary}

\begin{proof}
Given a computable sequence $(y_i)$, first compute the clamped sequence by
\cref{lem:clamp-computable}, then apply \cref{cor:G-sequential}.  By
\eqref{eq:totalized-inverse}, the result is $(\InvT(y_i))$.
\end{proof}

\subsection{Completion of the main theorem}

\begin{proof}[Proof of \cref{thm:main}]
Sequential computability is \cref{cor:G-sequential,cor:GT-sequential}.
Effective uniform continuity with the explicit recursive modulus is
\cref{thm:reciprocal-modulus}.  These are exactly the two clauses of
\cref{def:computable-function}.
\end{proof}

\subsection{An abstract inverse-computability principle}

The Fabius proof illustrates a reusable theorem.  Stating it clarifies which
properties are genuinely needed and which belong only to this example.

\begin{theorem}[Effective inversion on a computable compact interval]
\label{thm:abstract-inversion}
Let $f:[0,1]\to[0,1]$ be a strictly increasing homeomorphism.  Suppose:
\begin{enumerate}
\item $f(c)$ can be approximated uniformly and effectively at every dyadic
      rational $c$;
\item there is a computable positive rational sequence $(\alpha_p)$ such that
      for all $0\le x\le1-2^{-p}$,
      \begin{equation}\label{eq:abstract-gap}
       f(x+2^{-p})-f(x)\ge\alpha_p.
      \end{equation}
\end{enumerate}
Then $f^{-1}$ is sequentially computable and effectively uniformly continuous.
A tolerant-bisection algorithm is obtained by replacing $\Delta_r$ with
$\alpha_r$ in \cref{alg:tolerant-bisection}.
\end{theorem}

\begin{proof}
Condition \eqref{eq:abstract-gap} implies
$|u-v|<\alpha_p\Rightarrow|f^{-1}(u)-f^{-1}(v)|<2^{-p}$ by the same
contrapositive argument as \cref{cor:effective-injectivity}.  This gives a
binary recursive modulus.  The proof of
\cref{thm:bisection-correct} uses only this implication, dyadic evaluation of
$f$, and strict monotonicity, so it carries over verbatim.
\end{proof}

For the Fabius function, \cref{thm:least-mass} identifies the best possible
choice $\alpha_p=F(2^{-p})$; its minimizing intervals are exactly the left and
right endpoint intervals.  The box event in \cref{lem:box-event} supplies a
closed lower substitute requiring no invocation of the exact dyadic evaluator.

\section{Analytic consequences of the exact modulus}

\subsection{Uniform continuity without any H\"older exponent}

Effective uniform continuity is sometimes mentally associated with a power-law
modulus.  The inverse Fabius function is a clean counterexample to that
intuition.

\begin{theorem}[No positive H\"older exponent at the endpoints]
\label{thm:no-holder}
For every $\alpha>0$ and every $C>0$, there are arbitrarily small $y>0$ such
that
\begin{equation}\label{eq:no-holder-zero}
 G(y)>Cy^\alpha.
\end{equation}
By reflection, there are arbitrarily small $1-y>0$ such that
\begin{equation}\label{eq:no-holder-one}
 1-G(y)>C(1-y)^\alpha.
\end{equation}
Thus neither $G$ nor $\InvT$ is H\"older continuous of any positive order at
$0$ or $1$, despite being effectively uniformly continuous.
\end{theorem}

\begin{proof}
Choose an integer $q$ with $q\alpha>1$.  By
\cref{lem:flatness-upper}, for all sufficiently small $x>0$,
$F(x)\le C_qx^q$ with a fixed $C_q$.  If $G(y)\le Cy^\alpha$ held near zero,
substitution $y=F(x)$ would give
\[
 x=G(F(x))\le C F(x)^\alpha
 \le C C_q^\alpha x^{q\alpha}.
\]
After division by $x$, the right side tends to zero, a contradiction.  The
right endpoint follows from $G(1-y)=1-G(y)$.
\end{proof}

\begin{corollary}[The exact modulus is sub-H\"older]\label{cor:sub-holder}
For every $\alpha>0$,
\begin{equation}\label{eq:sub-holder}
 \frac{\omega_G(\eta)}{\eta^\alpha}
 =\frac{G(\eta)}{\eta^\alpha}\longrightarrow+\infty
 \qquad(\eta\downarrow0).
\end{equation}
\end{corollary}

\begin{proof}
Choose an integer $q$ with $q\alpha>1$ and write $x=G(\eta)$, so that
$\eta=F(x)$ and $x\downarrow0$ with $\eta\downarrow0$.  By
\cref{lem:flatness-upper}, for all sufficiently small $x$,
$F(x)\le C_qx^q$.  Hence
\[
 \frac{G(\eta)}{\eta^\alpha}
 =\frac{x}{F(x)^\alpha}
 \ge C_q^{-\alpha}x^{1-q\alpha}\longrightarrow+\infty.
\]
The equality with $\omega_G$ is \cref{thm:exact-inverse-modulus}.
\end{proof}

The endpoint growth used here is also formalized, independently of the new
exact-modulus module, in \path{FabiusInverseAsymptotic.lean}.  The declarations
\path{Fabius.rpow_isLittleO_fabiusInv_at_zero_right} and
\path{Fabius.tendsto_fabiusInv_div_rpow_atTop_at_zero_right} give the
sub-H\"older divergence, while
\path{Fabius.fabiusInv_not_isBigO_rpow_at_zero_right} and
\path{Fabius.not_exists_fabiusInv_le_const_mul_rpow_near_zero} give the
corresponding impossibility of a positive power-law upper bound.  These are
analytic inputs already present in the repository, not consequences claimed
to have been newly proved by \path{InverseModulus.lean}.

\subsection{Endpoint asymptotics become modulus asymptotics}

The live formal corpus proves the leading inverse equivalent through
\lean{fabiusInv_isEquivalent_fabiusInverseAsymptoticMain} in
\path{FabiusInverseAsymptotic.lean}.  The companion frontier reports give
stronger paper-level inverse expansions\ \cite{ProveIt}.
Let
\begin{equation}\label{eq:T-L-rho}
 T=\log(1/\eta),
 \qquad L=\log2,
 \qquad \rho=\sqrt{\frac{2T}{L}}.
\end{equation}
Suppressing the paper-level first log-periodic correction, the formalized leading
equivalent is
\begin{equation}\label{eq:G-leading}
 G(\eta)\sim
 G_0(\eta):=
 \frac{\sqrt T}{\e\sqrt L}
 \exp\!\bigl(-\sqrt{2LT}\bigr)
 \qquad(\eta\downarrow0).
\end{equation}
More precisely, \path{Inverse_Endpoint_All_Orders.tex} obtains at the
frontier-document level a relative expansion of the form
\begin{equation}\label{eq:G-periodic-correction}
 \frac{G(\eta)}{G_0(\eta)}
 =1+
 \frac{\frac12\log\rho+\kappa_0-\Psi(\theta)}{\rho}
 +\bigO\!\left(\frac{(\log\rho)^2}{\rho^2}\right),
\end{equation}
where $\Psi$ is a nonconstant smooth one-periodic Gamma--zeta fluctuation and
$\theta$ is the lower-Lambert phase shifted by an explicit constant.

\begin{corollary}[Paper-level transfer of the frontier inverse expansion]
\label{cor:sharp-asymptotic-modulus}
At the paper level, the expressions \eqref{eq:G-leading} and
\eqref{eq:G-periodic-correction} transfer, without any further inversion, to
the leading and first-correction asymptotics of the exact continuity modulus
of $G$ and $\InvT$:
\begin{equation}\label{eq:asymptotic-modulus}
 \omega_G(\eta)=\omega_{\InvT}(\eta)=G(\eta)
 \qquad(0\le\eta\le1).
\end{equation}
Likewise, every valid all-orders endpoint expansion for the inverse transfers
verbatim to an all-orders modulus-of-continuity expansion.  Lean currently
proves the leading equivalent and an implicit expansion pulled back to exact
inverse coordinates, but not the displayed explicit periodic correction or
an explicit all-orders reversion.
\end{corollary}

\begin{proof}
This is exactly \cref{thm:exact-inverse-modulus}.
\end{proof}

This observation closes a conceptual loop in the corpus.  The lower-Lambert
phase and Bell-polynomial inverse transfer were developed to approximate
quantiles.  The geometric theorem shows that the same quantiles are the
optimal global sensitivity profile of the inverse map.

\subsection{Derivative blow-up versus computability}

On $0<y<1$, the repository proves
\begin{equation}\label{eq:inverse-derivative}
 G'(y)=\frac{1}{F'(G(y))}
 =\frac{1}{2\Up(2G(y)-1)}>0.
\end{equation}
The derivative tends to $+\infty$ at both endpoints.  There is therefore no
global Lipschitz constant and no endpoint condition number bounded in the
ordinary differential sense.  Nonetheless the exact nonlinear modulus
$G(\eta)$ is computable, and tolerant bisection uses precisely that nonlinear
conditioning.  This is why derivative-based numerical intuition alone is not
a substitute for a computable-analysis proof.

\section{Numerical and exact-rational checks}

The accompanying script
\texttt{inverse\_fabius\_computability\_experiments.py} is a reproducibility
supplement.  The proofs do not depend on it.  It uses
\texttt{fractions.Fraction} for the Thue--Morse splines and asserts every
inequality it checks.

\subsection{Endpoint mass bounds}

Define
\begin{align}
 A_r&:=\frac{r(r-1)}2+\log_2(r!),\label{eq:A-r}\\
 B_r&:=\frac{r(r+1)}2+(r+1)\log_2(r+1).\label{eq:B-r}
\end{align}
Then \cref{lem:box-event,lem:flatness-upper} say
\begin{equation}\label{eq:AB-bracket}
 A_r\le\log_2\frac1{F(2^{-r})}\le B_r.
\end{equation}
The following values show the quadratic squeezing.

\begin{table}[htbp]
\centering
\small
\begin{tabular}{@{}r rr rr@{}}
\toprule
$r$ & $\Delta_r$ & $B_r=-\log_2\Delta_r$
    & flatness upper bound & $A_r$\\
\midrule
1  & $1.25000000\,10^{-1}$ & 3.000000
   & $1.00000000$          & 0.000000\\
2  & $4.62962963\,10^{-3}$ & 7.754888
   & $2.50000000\,10^{-1}$ & 2.000000\\
3  & $6.10351563\,10^{-5}$ & 14.000000
   & $2.08333333\,10^{-2}$ & 5.584963\\
4  & $3.12500000\,10^{-7}$ & 21.609640
   & $6.51041667\,10^{-4}$ & 10.584963\\
5  & $6.54097611\,10^{-10}$ & 30.509775
   & $8.13802083\,10^{-6}$ & 16.906891\\
6  & $5.79006996\,10^{-13}$ & 40.651484
   & $4.23855252\,10^{-8}$ & 24.491853\\
7  & $2.22044605\,10^{-16}$ & 52.000000
   & $9.46105473\,10^{-11}$ & 33.299208\\
8  & $3.75610368\,10^{-20}$ & 64.529325
   & $9.23931126\,10^{-14}$ & 43.299208\\
9  & $2.84217094\,10^{-24}$ & 78.219281
   & $4.01011773\,10^{-17}$ & 54.469133\\
10 & $9.72815993\,10^{-29}$ & 93.053748
   & $7.83226120\,10^{-21}$ & 66.791061\\
\bottomrule
\end{tabular}
\caption{Rigorous lower and upper bounds for $F(2^{-r})$.  The two logarithmic
columns bracket the optimal binary modulus index.}
\label{tab:endpoint-bounds}
\end{table}

The numerical width in \cref{tab:endpoint-bounds} is not a defect in the
computability proof.  Its order is only linear in $r$ after the common
$p^2/2+p\log_2p$ terms are removed.  Exact dyadic arithmetic or the sharp
Lambert expansion can replace the bracket whenever a tight practical modulus
is desired.

\subsection{Centered-spline convergence}

The exact values
\begin{equation}\label{eq:three-exact-values}
 F(1/2)=1/2,
 \qquad F(1/4)=5/72,
 \qquad F(1/8)=1/288
\end{equation}
provide compact checks of \eqref{eq:centered-spline-formula}.  The script
produces the following exact rational errors, displayed decimally.

\begin{table}[htbp]
\centering
\small
\begin{tabular}{@{}r c rr r@{}}
\toprule
$N$ & $x$ & $S_N(x)$ & $|S_N(x)-F(x)|$ & certified $2^{-N}$\\
\midrule
4  & $1/2$ & $5.0000000000000000\,10^{-1}$ & $0$ & $6.25\,10^{-2}$\\
4  & $1/4$ & $6.9010416666666671\,10^{-2}$ & $4.34027778\,10^{-4}$ & $6.25\,10^{-2}$\\
4  & $1/8$ & $3.2552083333333335\,10^{-3}$ & $2.17013889\,10^{-4}$ & $6.25\,10^{-2}$\\
\addlinespace
8  & $1/2$ & $5.0000000000000000\,10^{-1}$ & $0$ & $3.90625\,10^{-3}$\\
8  & $1/4$ & $6.9442749023437500\,10^{-2}$ & $1.69542101\,10^{-6}$ & $3.90625\,10^{-3}$\\
8  & $1/8$ & $3.4713745117187500\,10^{-3}$ & $8.47710503\,10^{-7}$ & $3.90625\,10^{-3}$\\
\addlinespace
12 & $1/2$ & $5.0000000000000000\,10^{-1}$ & $0$ & $2.44140625\,10^{-4}$\\
12 & $1/4$ & $6.9444437821706131\,10^{-2}$ & $6.62273831\,10^{-9}$ & $2.44140625\,10^{-4}$\\
12 & $1/8$ & $3.4722189108530679\,10^{-3}$ & $3.31136915\,10^{-9}$ & $2.44140625\,10^{-4}$\\
\bottomrule
\end{tabular}
\caption{Exact centered-spline checks.  The global certificate is intentionally
conservative; the pointwise dyadic errors are much smaller.}
\label{tab:spline-convergence}
\end{table}

For $N=8$, interval length $h=1/8$, and the dyadic grid $2^{-8}\Z$, the script
also verifies exactly that
\[
 \min_x\bigl(S_8(x+1/8)-S_8(x)\bigr)
 =S_8(1/8)=\frac{455}{131072},
\]
with minimizers $x=0$ and $x=7/8$.  This finite-convolution experiment
reproduces the complete continuous minimizer locus in \cref{thm:least-mass};
the theorem itself is the strict unimodality argument, not a grid inference.

\subsection{What the code deliberately does not claim}

The naive rational formula scales exponentially with $N$.  The script uses
small orders to check identities and error bounds, not to compete with the
repository's exact dyadic evaluator.  Likewise, its short bisection transcript
illustrates certified comparisons at a modest scale; sequential computability
is proved by the abstract finite algorithm, regardless of practical runtime.

\section{Connections to the wider Fabius--Rvachev corpus}
\label{sec:connections}

\subsection{Thue--Morse inclusion--exclusion}

The forward oracle is not an unrelated approximation scheme.  The binary
subset map in \cref{prop:finite-cdf} makes the Thue--Morse sign
$(-1)^{s_2(j)}$ the exact inclusion--exclusion coefficient of the finite
uniform convolution.  Thus the name-transformer for $G$ ultimately rests on
the same cancellation mechanism that drives the repository's exact dyadic
formulae and discrete-limit identities.

The centered shift $(2j+1)/2^{N+1}$ is equally structural: it replaces the
unknown tail by its mean and halves the elementary uniform error from the
uncentered bound $2^{1-N}$ to $2^{-N}$.  In computational terms, this one-bit
improvement is free.

\subsection{\texorpdfstring{$q$}{q}-binomial coefficients and accelerated forward oracles}

The finite Thue--Morse sums admit Gaussian-binomial reorganizations; the
arithmetic background is developed in \cite{Arias06487,ProveIt}.  The
inverse/sampling corpus uses $q=1/4$ Richardson filters on the geometric error
scale of centered finite approximants; the current snapshot also adds a
quarter-base $q$-Pochhammer factorization of the Rvachev sinc product.  None of
these identities is needed for mere computability, because
\cref{thm:centered-error} already gives an effective rate.  They are relevant
to complexity: a certified accelerated approximation to $F(c)$ can be dropped
into \cref{alg:tolerant-bisection} without changing one line of its logical
proof, provided the accelerated remainder is enclosed rigorously.

\subsection{Infinite sinc products and alternative evaluation}

Equation \eqref{eq:sinc-product} gives the Fourier transform of $\Up$, in
the classical infinite-convolution tradition of
\cite{JessenWintner1935,Arias05442}.  Fourier inversion therefore offers an independent evaluator for $F$ or its density.
The corpus develops strong decay estimates, Euler/canonical products, and
Poisson-summation formulae.  A computability proof based on this route would
need an explicit truncation error for both the frequency integral and the
infinite product.  Once those are supplied, tolerant bisection is again
unchanged.

This modularity is useful for cross-certification: finite splines work in
physical space with exact rationals; sinc products work in frequency space
with analytic tail bounds.  Agreement of independently certified enclosures
is a stronger implementation check than agreement of unproved decimal
approximations.

\subsection{Lambert \texorpdfstring{$W$}{W}, Bell polynomials, and Bernoulli data}

The flat endpoint makes the inverse poorly conditioned in any power scale.
The lower branch $W_{-1}$ naturally solves the saddle equation controlling
$F(x)$ as $x\downarrow0$.  The frontier reports then use Bell-polynomial
reversion to transport all-orders forward expansions to $G(y)$ \cite{ProveIt}; Bernoulli numbers and
polynomials enter through logarithms of the Laplace/sinc products and through
formal composition.

At the paper level, the exact identity $\omega_G(\eta)=G(\eta)$ gives these
expansions a second interpretation: they are not only quantile asymptotics but
optimal continuity-modulus asymptotics.  In particular, the nonconstant
periodic Gamma--zeta term cannot be discarded from a genuinely sharp
effective modulus.  This interpretation is not yet a public Lean theorem.

\subsection{Legendre and Lagrange representations}

The corpus expands $\Up$ in even Legendre polynomials and represents
polynomials by finite combinations of shifted atomic functions.  Such bases
can provide spectral or interpolation-based forward oracles.  For a formal
computability theorem, however, convergence alone is insufficient: one needs
a recursive tail majorant uniform in the evaluation point.  Deriving sharp
certified Legendre-tail bounds and comparing their bit complexity with the
Thue--Morse spline is a natural next step.

\subsection{Total positivity, unimodality, and interval masses}

The least-mass lower bound uses only symmetric unimodality; its exact
endpoint-only equality locus uses strict radial unimodality.  Several corpus
reports investigate still stronger total-positivity and variation-diminishing
properties.  If a stronger kernel theorem is established, it may control not
only one interval mass but higher finite differences and multidimensional
quantile maps.  The computability proof therefore isolates a low-order shadow
of a potentially broader shape theory.

\section{Lean formalization status and remaining roadmap}
\label{sec:lean-status}

\subsection{Why the present \texttt{noncomputable} definition is not an obstacle}

The repository defines \texttt{fabiusInv} by inverting an order isomorphism and
then applying \texttt{IccExtend}.  Lean marks that definition
\texttt{noncomputable}, because the inverse equivalence is constructed by
classical existence machinery.  This is a statement about the chosen term in
Lean's programming language, not a theorem that the represented real function
has no algorithm.

Computable analysis deliberately separates these two levels.  An extensional
real function may be introduced by a classical characterization and later be
shown to possess a computable \emph{realizer}: a program transforming names of
inputs into names of outputs.  \Cref{alg:tolerant-bisection} is exactly such a
realizer.  Its correctness is proved by showing that its output denotes the
same extensional function as \texttt{fabiusInv}.

This distinction already occurs on the forward side.  The analytic canonical
Fabius function is defined abstractly, whereas
\path{FabiusComputableSpline.lean} supplies a primitive-recursive evaluator
and proves that it realizes the same function.  The inverse formalization can
follow the same architecture.

\subsection{The completed structural inverse-modulus layer}

The structural part of \cref{sec:exact-modulus} now lives in
\begin{center}
\path{Analysis/FabiusFunction/Lean/FabiusFunction/InverseModulus.lean}.
\end{center}
It imports \path{FabiusInverse.lean} and deliberately remains independent of
the heavier random-series and computable-analysis interfaces.  All thirty-two
public declarations (one definition and thirty-one theorems) were
compiler-checked as one focused target; each theorem uses only the standard
logical principles already present in its dependencies.  The declaration-level
correspondence is as follows.

\begin{longtable}{@{}>{\raggedright\arraybackslash}p{0.31\textwidth}>{\raggedright\arraybackslash}p{0.60\textwidth}@{}}
\toprule
Report object or claim & Lean declaration \\
\midrule
\endhead
The fixed-length increment $M_h(x)=F(x+h)-F(x)$ &
\path{Fabius.fabiusIntervalMass}.\\
\addlinespace
Reflection $M_h(1-h-x)=M_h(x)$ &
\path{Fabius.fabiusIntervalMass_reflect}.\\
\addlinespace
The zero tails $M_h(x)=0$ when $x+h\le0$ or $1\le x$, for $h\ge0$ &
\path{Fabius.fabiusIntervalMass_eq_zero_of_add_nonpos} and
\path{Fabius.fabiusIntervalMass_eq_zero_of_one_le}.\\
\addlinespace
Global left-half monotonicity and right-half antitonicity &
\path{Fabius.monotoneOn_fabiusIntervalMass_firstHalf} and
\path{Fabius.antitoneOn_fabiusIntervalMass_secondHalf}.  These are the weak
half-line statements that include the clamped zero tails.\\
\addlinespace
Maximal strict branches $[-h,(1-h)/2]$ and $[(1-h)/2,1]$ for $0<h\le1$ &
\path{Fabius.strictMonoOn_fabiusIntervalMass_firstHalf} and
\path{Fabius.strictAntiOn_fabiusIntervalMass_secondHalf}.\\
\addlinespace
Least endpoint increment $F(b-a)\le F(b)-F(a)$ &
\path{Fabius.fabiusReal_sub_le_sub}.\\
\addlinespace
Strict interior increment $F(h)<M_h(x)$ and exact endpoint starts &
\path{Fabius.fabiusReal_lt_fabiusIntervalMass_of_mem_Ioo} and
\path{Fabius.fabiusIntervalMass_eq_fabiusReal_iff}.\\
\addlinespace
Strict interior interval inequality and complete equality law
\eqref{eq:least-mass-equality-ab} &
\path{Fabius.fabiusReal_sub_lt_sub} and
\path{Fabius.fabiusReal_sub_eq_sub_iff}.\\
\addlinespace
Constrained forward superadditivity \eqref{eq:fabius-superadditive} &
\path{Fabius.fabiusReal_add_le}, with equality classification
\eqref{eq:fabius-superadditive-equality} in
\path{Fabius.fabiusReal_add_eq_iff}.\\
\addlinespace
Ordered inverse gap on $[0,1]$ &
\path{Fabius.fabiusInv_sub_le_sub_of_mem_Icc}, with exact equality cases
\eqref{eq:inverse-gap-equality} in
\path{Fabius.fabiusInv_sub_eq_sub_iff_of_mem_Icc}.\\
\addlinespace
Ordered inverse gap including clamped tails, and its order-free strengthening
\eqref{eq:inverse-gap-bound} &
\path{Fabius.fabiusInv_sub_le_sub_of_le} and
\path{Fabius.fabiusInv_sub_le_sub}.\\
\addlinespace
Global inverse subadditivity \eqref{eq:inverse-subadditive} &
\path{Fabius.fabiusInv_add_le}.\\
\addlinespace
Absolute and metric pointwise modulus \eqref{eq:absolute-inverse-gap-bound} &
\path{Fabius.abs_fabiusInv_sub_le} and
\path{Fabius.dist_fabiusInv_le}; the complete unit-input equality locus
\eqref{eq:absolute-inverse-gap-equality} is
\path{Fabius.abs_fabiusInv_sub_eq_iff_of_mem_Icc}.\\
\addlinespace
The saturation identity $\InvT(\min(\eta,1))=\InvT(\eta)$ &
\path{Fabius.fabiusInv_min_one}.\\
\addlinespace
Attained exact totalized modulus and its supremum identity &
\path{Fabius.isGreatest_abs_fabiusInv_sub} and
\path{Fabius.sSup_abs_fabiusInv_sub_eq}.\\
\addlinespace
Attained exact modulus with unit-interval inputs, for every $\eta\ge0$ &
\path{Fabius.isGreatest_abs_fabiusInv_sub_Icc} and
\path{Fabius.sSup_abs_fabiusInv_sub_Icc_eq}.\\
\addlinespace
Strict and closed effective injectivity &
\path{Fabius.abs_fabiusInv_sub_lt_of_abs_sub_lt_fabiusReal} and
\path{Fabius.abs_fabiusInv_sub_le_of_abs_sub_le_fabiusReal}.\\
\addlinespace
Contrapositive separation \eqref{eq:contrapositive-effective-injectivity} &
\path{Fabius.fabiusReal_le_abs_sub_of_le_abs_fabiusInv_sub}.\\
\addlinespace
Exact strict-threshold biconditional \eqref{eq:exact-strict-threshold} &
\path{Fabius.forall_abs_fabiusInv_sub_lt_iff}.\\
\bottomrule
\end{longtable}

The formal least-mass API deliberately separates four reusable layers: the
global weak shape, the maximal strict branches, the lower bound, and the
complete equality locus.  The inverse equality theorems remain restricted to
unit inputs because the clamped tails add straddling-tail equality cases and
make the endpoint alternatives order-dependent globally.  The
\texttt{IsGreatest} declarations record both the sharp upper bound and an
extremizing pair; the \texttt{sSup} equalities are immediate corollaries.  The
totalized formal theorem returns \texttt{fabiusInv ... eta}, which is literally
the report's $G(\min(\eta,1))$ because \texttt{fabiusInv} is already clamped.

\subsection{The completed effective-uniform-continuity layer}

The next formal layer is
\begin{center}
\path{Analysis/FabiusFunction/Lean/FabiusFunction/FabiusInverseEffectiveContinuity.lean}.
\end{center}
It imports the forward computability interface, the exact inverse-dyadic
recurrence, and \path{InverseModulus.lean}.  Its recurrence proof is stronger
than the report's box estimate, while its final reciprocal modulus is simpler
and larger than the logarithmic-$r(n)$ modulus.  The following table exhausts
all fourteen public declarations in the module (two definitions and twelve
theorems); private recurrence and
primitive-recursion helpers are intentionally omitted.

\begin{longtable}{@{}>{\raggedright\arraybackslash}p{0.34\textwidth}>{\raggedright\arraybackslash}p{0.57\textwidth}@{}}
\toprule
Lean declaration & Report object or claim \\
\midrule
\endhead
\path{Fabius.inverseFabiusFactorialDenominator} &
The recursively executable factorial denominator $d_{\mathrm{fac}}(r)$.\\
\addlinespace
\path{Fabius.inverseFabiusFactorialDenominator_eq} &
The closed form
$d_{\mathrm{fac}}(r)=2^{\binom{r+1}{2}}(r+1)!$.\\
\addlinespace
\path{Fabius.inverseFabiusFactorialDenominator_primrec} &
Primitive recursiveness of $d_{\mathrm{fac}}$.\\
\addlinespace
\path{Fabius.inverseFabiusDeltaDenominator} &
The report's elementary denominator
$d_{\Delta}(r)=2^{\binom{r+1}{2}}(r+1)^{r+1}=\Delta_r^{-1}$.\\
\addlinespace
\path{Fabius.inverseFabiusDeltaDenominator_primrec} &
Primitive recursiveness of $d_{\Delta}$.\\
\addlinespace
\path{Fabius.fabiusReal_inverse_two_pow_one_term_lower_bound} &
For $r>0$, the sharper unrounded zeroth-term recurrence bound with denominator
$2^{\binom r2}(r+1)!(2^r-1)$.\\
\addlinespace
\path{Fabius.inv_inverseFabiusFactorialDenominator_le_fabiusReal} &
The first inequality in \eqref{eq:factorial-endpoint-lower}, including $r=0$.\\
\addlinespace
\path{Fabius.inverseFabiusFactorialDenominator_le_deltaDenominator} &
The denominator comparison $d_{\mathrm{fac}}(r)\le d_{\Delta}(r)$, from
$(r+1)!\le(r+1)^{r+1}$.\\
\addlinespace
\path{Fabius.inv_inverseFabiusDeltaDenominator_le_fabiusReal} &
The exact numerical conclusion $\Delta_r\le F(2^{-r})$ of
\cref{lem:box-event}, proved through the recurrence rather than its box event.\\
\addlinespace
\path{Fabius.abs_fabiusInv_sub_lt_inverse_two_pow_of_lt_factorialDenominator} &
The stronger strict dyadic inverse modulus: input distance
$<d_{\mathrm{fac}}(r)^{-1}$ forces output distance $<2^{-r}$.\\
\addlinespace
\path{Fabius.abs_fabiusInv_sub_le_inverse_two_pow_of_le_factorialDenominator} &
The closed factorial companion: input distance
$\le d_{\mathrm{fac}}(r)^{-1}$ forces output distance $\le2^{-r}$.\\
\addlinespace
\path{Fabius.abs_fabiusInv_sub_lt_inverse_two_pow_of_lt_deltaDenominator} &
The report's exact strict implication \eqref{eq:closed-dyadic-modulus}: input
distance $<\Delta_r$ forces output distance $<2^{-r}$.\\
\addlinespace
\path{Fabius.abs_fabiusInv_sub_le_inverse_two_pow_of_le_deltaDenominator} &
The exact closed-boundary companion: input distance $\le\Delta_r$ forces
output distance $\le2^{-r}$.\\
\addlinespace
\path{Fabius.fabiusInv_effectivelyUniformContinuous} &
Effective uniform continuity of the totalized inverse, witnessed by
$n\mapsto d_{\mathrm{fac}}(n)$ and hence using the simple scale $r=n$.\\
\bottomrule
\end{longtable}

This table crosswalks theorem statements, not proof methods.  In particular,
the declaration with conclusion $\Delta_r\le F(2^{-r})$ does not construct the
random-series event used in \cref{lem:box-event}.  Conversely, the final
theorem establishes the repository predicate
\texttt{EffectivelyUniformContinuous}; it does not establish a sequential
realizer or the two-field \texttt{IsComputableRealFunction} structure.

\subsection{The completed logarithmic reciprocal-modulus layer}

The smaller report-facing refinement is
\begin{center}
\path{Analysis/FabiusFunction/Lean/FabiusFunction/FabiusInverseLogarithmicModulus.lean}.
\end{center}
It imports the preceding effective-continuity module, constructs the least
dyadic order by primitive-recursive bounded search, and composes both explicit
endpoint-mass denominators with that order.  The following table exhausts all
eighteen public declarations in the module (three definitions and fifteen
theorems); private power, search-relation, and real reciprocal-comparison
helpers are intentionally omitted.

\begin{longtable}{@{}>{\raggedright\arraybackslash}p{0.39\textwidth}>{\raggedright\arraybackslash}p{0.52\textwidth}@{}}
\toprule
Lean declaration & Report object or claim \\
\midrule
\endhead
\path{Fabius.inverseFabiusLogarithmicOrder} &
The bounded-search selector $r(n)$, totalized by the harmless convention
$r(0)=1$.\\
\addlinespace
\path{Fabius.inverseFabiusLogarithmicOrder_eq_succ_log2} &
The binary-length identity $r(n)=\operatorname{log}_2(n)+1$, with
\texttt{Nat.log2} understood in its natural-valued sense.\\
\addlinespace
\path{Fabius.inverseFabiusLogarithmicOrder_primrec} &
Primitive recursiveness of the bounded-search selector.\\
\addlinespace
\path{Fabius.inverseFabiusLogarithmicOrder_isLeast} &
For $n>0$, the exact least-order characterization
$r(n)=\min\{r\in\N:n<2^r\}$ from \cref{thm:reciprocal-modulus}.\\
\addlinespace
\path{Fabius.inverseFabiusLogarithmicOrder_le_self} &
For $n>0$, the logarithmic order does not exceed the former coarse choice
$r=n$.\\
\addlinespace
\path{Fabius.inverseFabiusLogarithmicFactorialDenominator} &
The stronger composed denominator $d_{\mathrm{fac}}(r(n))$ for $n>0$, with
value $1$ at zero.\\
\addlinespace
\path{Fabius.inverseFabiusLogarithmicDeltaDenominator} &
The report-exact denominator
$d(n)=d_{\Delta}(r(n))=2^{r(n)(r(n)+1)/2}(r(n)+1)^{r(n)+1}$ for $n>0$, with
$d(0)=1$.\\
\addlinespace
\path{Fabius.inverseFabiusLogarithmicFactorialDenominator_of_pos} &
The positive-input unfolding to $d_{\mathrm{fac}}(r(n))$.\\
\addlinespace
\path{Fabius.inverseFabiusLogarithmicDeltaDenominator_of_pos} &
The positive-input unfolding to the report's $d_{\Delta}(r(n))$.\\
\addlinespace
\path{Fabius.inverseFabiusLogarithmicFactorialDenominator_primrec} &
Primitive recursiveness of the stronger composed denominator.\\
\addlinespace
\path{Fabius.inverseFabiusLogarithmicDeltaDenominator_primrec} &
Primitive recursiveness of the report-exact $d(n)$.\\
\addlinespace
\path{Fabius.inverseFabiusLogarithmicFactorialDenominator_le_deltaDenominator} &
The pointwise comparison between the stronger factorial and report-exact
Delta denominators.\\
\addlinespace
\path{Fabius.abs_fabiusInv_sub_lt_inv_nat_of_lt_logarithmicFactorialDenominator} &
The stronger reciprocal modulus with strict input threshold: input distance
$<d_{\mathrm{fac}}(r(n))^{-1}$ forces output distance $<1/n$.\\
\addlinespace
\path{Fabius.abs_fabiusInv_sub_lt_inv_nat_of_le_logarithmicFactorialDenominator} &
Its closed-input-threshold companion; the output remains strictly below
$1/n$ because $2^{-r(n)}<1/n$.\\
\addlinespace
\path{Fabius.abs_fabiusInv_sub_lt_inv_nat_of_lt_logarithmicDeltaDenominator} &
The strict-input form of the exact implication claimed in
\cref{thm:reciprocal-modulus}, using the report's $d(n)$.\\
\addlinespace
\path{Fabius.abs_fabiusInv_sub_lt_inv_nat_of_le_logarithmicDeltaDenominator} &
The closed-input-threshold strengthening for the same $d(n)$; its output
conclusion is still strict.\\
\addlinespace
\path{Fabius.fabiusInv_effectivelyUniformContinuous_logarithmic} &
\texttt{EffectivelyUniformContinuous} for the totalized inverse, witnessed by
the stronger logarithmic factorial denominator.\\
\addlinespace
\path{Fabius.fabiusInv_effectivelyUniformContinuous_logarithmicDelta} &
The same repository predicate packaged with the report-exact logarithmic
Delta denominator.\\
\bottomrule
\end{longtable}

Together, the least-order, Delta-denominator, primitive-recursion, and inverse
modulus declarations are an exact formal counterpart of
\cref{thm:reciprocal-modulus}; the factorial declarations provide a stronger
alternative at the same order.  This module does not formalize the
probabilistic box event used in the human proof of the endpoint lower bound,
the exact ceiling $d_*$, tolerant bisection, sequential computability, the
combined computable-real-function theorem, or input-bit asymptotics.

\subsection{The remaining name-transformation module}

A natural downstream file remains
\begin{center}
\path{Analysis/FabiusFunction/Lean/FabiusFunction/}\\
\path{FabiusInverseComputability.lean}.
\end{center}
It should import
\path{FabiusFunction.FabiusInverseLogarithmicModulus}.  It no longer needs to
reprove interval shape, inverse-gap estimates, dyadic endpoint lower bounds,
or effective uniform continuity.  Its remaining public surface is the
name-transformation layer; the following Lean is schematic rather than a
claim about declarations already present.

\begin{lstlisting}[style=reportpseudo,language={}]
/-- Tolerant bisection transforms every computable input sequence. -/
theorem fabiusInv_sequentiallyComputable :
    SequentiallyComputable (fabiusInv fabius fabius_spec)

/-- Both Grzegorczyk clauses, packaged together. -/
theorem fabiusInv_isComputableRealFunction :
    IsComputableRealFunction (fabiusInv fabius fabius_spec)
\end{lstlisting}

\subsection{Remaining formalization order}

The former structural roadmap steps are complete in
\path{InverseModulus.lean}, and the closed-bound/effective-continuity steps are
complete in \path{FabiusInverseEffectiveContinuity.lean}.  The logarithmic
selector and report denominator are complete in
\path{FabiusInverseLogarithmicModulus.lean}.  The remaining order is:

\begin{enumerate}
\item If a probability-level counterpart is desired, formalize the particular
      box event in \eqref{eq:box-restrictions}, its containment, and its product
      probability.  This is no longer needed to obtain its endpoint
      inequality, which is already formalized by recurrence.
\item Connect the exact dyadic evaluator to the positive rational
      $F(2^{-r})$ and package the reciprocal-ceiling modulus, including the
      qualified fixed-dyadic-target minimality statement of
      \cref{prop:exact-dyadic-modulus}.
\item Formalize a generic tolerant-bisection lemma for a monotone computable
      map equipped with a recursive lower-increment bound.  The result should
      not be Fabius-specific.
\item Instantiate the generic lemma with the primitive-recursive centered
      spline evaluator from \path{FabiusComputableSpline.lean}.
\item Package the resulting sequential-computability theorem together with
      the now-formalized effective uniform continuity into the repository's
      computable-real-function interface.
\item Only after the effective core is complete, formalize the flatness bounds
      and input-bit asymptotics; they are not logical prerequisites of the
      computability theorem.
\end{enumerate}

\subsection{Reusable generic infrastructure}

The abstract \cref{thm:abstract-inversion} deserves its own module.  In a
representation-oriented library it can be stated using an oracle for $f$ and a
recursive lower-separation function
\[
  \alpha(p)>0,
  \qquad x+2^{-p}\le b
  \Longrightarrow f(x+2^{-p})-f(x)\ge\alpha(p).
\]
The conclusion is a computable realizer for the inverse.  Tolerant bisection
then becomes library infrastructure usable for distribution quantiles,
implicit monotone roots, cumulative spectral measures, and monotone solutions
of integral equations.  The proof is constructive in the strong sense that
all precision budgets are explicit.

\section{Further deductions, conjectures, and research directions}
\label{sec:frontiers}

\subsection{A certified practical inverse evaluator}

The proof algorithm is intentionally elementary rather than efficient.  A
practical certified implementation should retain its logical shell while
replacing the forward oracle by a portfolio of enclosures:

\begin{itemize}
\item exact dyadic evaluation and terminating local Taylor formulae near
      dyadic anchors;
\item centered Thue--Morse splines at moderate depth;
\item quarter-base $q$-Richardson acceleration of finite-prefix errors;
\item Fourier/sinc evaluation in regions where frequency-space decay is
      favorable;
\item lower-Lambert endpoint enclosures when the requested target is extremely
      close to $0$ or $1$.
\end{itemize}

Each backend must return a rational interval containing $F(c)$, not merely a
floating-point approximation.  The bisection controller only asks whether
that interval lies wholly above or below the target interval; overlapping
intervals trigger refinement or the modulus-based acceptance branch.  Thus
backend selection is a performance concern, not part of the correctness
argument.

\begin{problem}[Near-optimal certified evaluation]
Construct and formally verify an inverse evaluator whose oracle-bit
consumption is $\mu(p)+\bigO(\log p)$ for $p$ output bits at the worst endpoint,
and whose arithmetic running time is quasi-linear or softly quadratic in the
size of the exact integers produced.  Separate information-theoretic oracle
precision from the cost of evaluating $F$ to that precision.
\end{problem}

\subsection{Sharp integer moduli and the log-periodic phase}

The real-valued optimal input-bit threshold is
\[
 \Lambda(p):=\log_2\frac1{F(2^{-p})},
 \qquad \mu(p)=\ceil{\Lambda(p)}.
\]
The all-orders endpoint theory expresses $\Lambda(p)$ in a lower-Lambert
phase with a nonconstant one-periodic Gamma--zeta contribution and a hierarchy
of inverse powers.  Passing to the integer ceiling introduces a second,
sawtooth nonlinearity.  It is therefore natural to study the arithmetic
interaction between the smooth periodic saddle phase and fractional parts of
its growing nonperiodic carrier.

\begin{conjecture}[Nontrivial fluctuation of the optimal integer modulus]
\label{conj:integer-modulus}
After subtracting any fixed finite truncation of the full nonperiodic
asymptotic expansion of $\Lambda(p)$, the bounded residual governing
$\mu(p)$ does not become eventually constant or eventually periodic in the
integer variable $p$.  Its accumulation set is generated jointly by the
Gamma--zeta phase and the ceiling sawtooth.
\end{conjecture}

The conjecture is deliberately qualitative.  A stronger equidistribution
claim would require information about fractional parts of a nonlinear
Lambert-type phase that is not supplied by the present corpus.  Exact values
$F(2^{-p})\in\Q$ make the sequence amenable to arbitrary-precision
experimentation without numerical ambiguity.

\subsection{Base-\texorpdfstring{$b$}{b} atomic distributions}

Replace the dyadic random series by
\[
 X_b=\sum_{k\ge1}b^{-k}U_k,
 \qquad b\ge2,
\]
with independent uniform coordinates.  The finite convolution is still
computable, but binary Thue--Morse signs are replaced by base-$b$ subset or
digit combinatorics, and the support is $[0,1/(b-1)]$ before normalization.
The least-interval-mass theorem continues to follow whenever the normalized
density is symmetric and unimodal.  The same proof would then identify the
optimal inverse modulus with the quantile itself.

\begin{problem}[Uniform computability across the base]
Develop a theorem uniform in the integer parameter $b$: construct a single
algorithm that, given $b$, a name of $y$, and a precision $p$, computes the
normalized quantile of $X_b$.  Determine the leading dependence of the
optimal input-bit modulus on both $b$ and $p$, and identify the corresponding
base-$b$ log-periodic correction.
\end{problem}

\subsection{Certified spectral and polynomial representations}

The Legendre, Lagrange, and shifted-$\Up$ representations in the corpus are
mathematically exact, but a representation becomes an algorithm only after a
recursive uniform tail bound is attached.  This suggests three concrete
projects:

\begin{enumerate}
\item derive explicit Legendre coefficient envelopes strong enough to certify
      pointwise and uniform evaluation of $\Up$;
\item combine the finite shifted-$\Up$ representation of polynomials with
      interval arithmetic to obtain certified local surrogates for $F$;
\item compare physical-space spline, Legendre, and sinc-product evaluators by
      bit complexity, not merely by observed decimal convergence.
\end{enumerate}

A cross-certified implementation could evaluate the same point by two
independent representations and intersect the resulting enclosures.  This is
particularly attractive near transition regions where one representation's
natural error bound is pessimistic.

\subsection{Higher-dimensional monotone transport}

The one-dimensional theorem
\[
 \min_x \Prob\{X\in(x,x+h]\}=\Prob\{X\in(0,h]\}
\]
is a sharp anti-concentration lower bound tailored to a symmetric unimodal
law.  Because the law is atomless, the same identity holds with any choice of
open or closed endpoints.  In higher dimensions, inverse CDFs are replaced by
triangular Rosenblatt maps or optimal-transport maps.  Effective invertibility
would require computable lower bounds on masses of suitable boxes or sections.

\begin{problem}[Atomic transport computability]
For products or coupled variants of Fabius-type laws, determine conditions
under which the triangular transport from Lebesgue measure has a computable
inverse with an explicit multivariate modulus.  Investigate whether total
positivity or variation-diminishing properties supply the needed lower
section-mass estimates.
\end{problem}

\subsection{Branchwise inversion of the signed global extension}

The signed global Fabius extension is not globally injective, so it has no
single inverse.  On every maximal monotonicity interval, however, it has a
branch inverse.  Local finiteness of its Thue--Morse translate expansion makes
forward evaluation computable.  The remaining issue is effective branch
separation near critical values.

\begin{problem}[Computable branch atlas]
Construct an effective enumeration of monotonicity intervals of the signed
global extension, together with certified branch ranges and recursive inverse
moduli.  Determine whether the self-similar translation law permits one finite
branch template plus a Thue--Morse address.
\end{problem}

\section{Conclusion}
\label{sec:conclusion}

At the paper level, the inverse Fabius function is computable in precisely the
two-part sense requested.  The effective-uniform-continuity clause follows
from the exact
geometric identity
\[
  \omega_G(\eta)=G(\min(\eta,1)),
\]
combined with the closed positive lower bound
\[
 F(2^{-r})\ge
 2^{-r(r+1)/2}(r+1)^{-(r+1)}.
\]
The sequential-computability clause follows from a tolerant bisection
transformer that uses only finite exact rational operations, a computable name
of the target, and certified centered-spline evaluations of $F$.  It never
requires deciding equality of computable reals.

The result is stronger than an existence proof.  It identifies the
\emph{optimal} modulus, explains why inversion loses quadratically many bits
at the flat endpoints, and turns the frontier lower-Lambert quantile
asymptotics into sharp paper-level continuity-modulus asymptotics.  The same proof covers
the inverse on $[0,1]$ and the canonical clamped totalization on all of
$\R$.

Conceptually, the decisive bridge is simple: weak symmetric unimodality makes
endpoint intervals minimize mass, while strict radial unimodality identifies
every equality case.  In the Fabius setting that sharpened shape fact connects
probability, computable analysis, exact dyadic arithmetic, and inverse
asymptotics.  It is also the part most likely to generalize beyond the
particular atomic function studied here.

The current Lean boundary cuts through the two clauses of computability.
\path{InverseModulus.lean} certifies the exact structural modulus, and
\path{FabiusInverseEffectiveContinuity.lean} certifies the closed dyadic bound
and \texttt{EffectivelyUniformContinuous}, using a recurrence estimate that is
stronger than the report's box bound.
\path{FabiusInverseLogarithmicModulus.lean} certifies the least
logarithmic-$r(n)$ selector and both the report-exact and stronger factorial
reciprocal moduli.  The probabilistic box-event proof, exact-ceiling modulus,
tolerant bisection, sequential computability, final computable-real-function
package, and precision asymptotics remain to be formalized.

\appendix

\section{Pseudocode with explicit error budgets}
\label{app:pseudocode}

The following version makes every tolerance in
\cref{alg:tolerant-bisection} explicit.  The function
\texttt{TargetName} returns, on inputs $i,s$, a rational within $2^{-s}$ of
the clamped input $y_i$.  The function \texttt{CenteredSpline} returns, on
inputs $N,c$, the exact rational
$S_N(c)$ of \eqref{eq:centered-spline-formula}.

\begin{lstlisting}[style=reportpseudo,language={}]
InverseFabiusName(i, p):
    # Requested absolute output error: 2^(-p).
    r     := p + 2
    h     := 2^(-r)
    delta := 2^(-r*(r+1)/2) * (r+1)^(-(r+1))

    choose s with 2^(-s) <= delta/8
    choose N with 2^(-N) <= delta/8
    yhat := clamp(TargetName(i, s), 0, 1)

    a := 0
    b := 1

    repeat p + 2 times:
        c    := (a + b)/2
        fhat := CenteredSpline(N, c)
        D    := fhat - yhat

        # |D - (F(c)-y_i)| <= delta/4.
        if D > delta/2:
            b := c                # certified F(c) > y_i
        else if D < -delta/2:
            a := c                # certified F(c) < y_i
        else:
            return c              # |F(c)-y_i| < delta

    return (a + b)/2
\end{lstlisting}

In the acceptance branch,
$\abs{F(c)-y_i}<3\delta/4<\delta\le F(h)$, so
\cref{cor:effective-injectivity} gives
$\abs{c-G(y_i)}<h<2^{-p}$.  If no acceptance occurs, every update is
certified and the true inverse remains in $[a,b]$; after $p+2$ iterations its
width is $2^{-(p+2)}$, so the final midpoint is within
$2^{-(p+3)}<2^{-p}$.

The constants have generous slack.  They were chosen to make the proof
transparent and to absorb representation-dependent rounding.  Replacing the
three occurrences of $1/8,1/4,1/2$ by any fixed compatible margin gives the
same computability theorem.

\section{Proof-dependency and status audit}
\label{app:dependency-audit}

\begin{longtable}{@{}>{\raggedright\arraybackslash}p{0.28\textwidth}>{\raggedright\arraybackslash}p{0.32\textwidth}>{\raggedright\arraybackslash}p{0.31\textwidth}@{}}
\toprule
Claim & Input & Human proof and Lean status \\
\midrule
\endhead
Random-series representation of $F$ & Infinite convolution / probability
representation in the repository and classical Fabius theory & Imported;
normalization restated in \cref{sec:forward}.\\
\addlinespace
Exact finite CDF \eqref{eq:finite-cdf-formula} & Inclusion--exclusion for finite sums
of uniforms & Paper-level proof in \cref{prop:finite-cdf}.\\
\addlinespace
Uniform centered-spline error $2^{-N}$ & Tail support and the global density
bound $p\le2$ & Paper-level proof in \cref{thm:centered-error}; formalized
counterparts include
\lean{abs_uniformCenteredPartialCDF_sub_fabiusReal_le} and
\lean{abs_fabiusUniformSpline_sub_fabiusReal_le_all}.\\
\addlinespace
Strict increase of $F$ & Positivity of finite-convolution densities and a
short tail & Paper-level proof in \cref{prop:strict-increase}; Lean proves
\lean{strictMonoOn_fabiusReal}.\\
\addlinespace
Symmetric strict unimodality of $p=F'$ & Differential equation, symmetry, and
strict increase of $F$ & Paper-level derivation in
\cref{prop:density-shape}; Lean proves
\lean{strictMonoOn_deriv_fabiusReal_Icc} and
\lean{strictAntiOn_deriv_fabiusReal_Icc}.\\
\addlinespace
Least fixed-length increment $F(h)$ & Symmetric unimodality of $p$ & Human
proof in \cref{thm:least-mass}; reflection, global weak shape, maximal strict
branches, the lower bound, and the complete degenerate/endpoint equality locus
are formalized in \path{InverseModulus.lean}.  Exact names appear in
\cref{sec:lean-status}.\\
\addlinespace
Exact inverse modulus $G(\eta)$ & Least interval mass and inverse identities &
Human proof in \cref{thm:exact-inverse-modulus}; attained maxima and exact
suprema are formalized in \path{InverseModulus.lean}.  Exact names appear in
\cref{sec:lean-status}.\\
\addlinespace
Exact effective-injectivity threshold $F(h)$ & Pointwise inverse modulus and
inverse identities & Human proofs in
\cref{cor:effective-injectivity,thm:exact-strict-threshold}; strict, closed,
contrapositive, and biconditional forms are formalized in
\path{InverseModulus.lean}.\\
\addlinespace
Closed modulus $\Delta_r$ & Explicit positive box event in the random series &
The event proof is in \cref{lem:box-event}; its numerical endpoint inequality
and the exact inverse implication of \cref{thm:closed-dyadic-modulus} are
formalized in \path{FabiusInverseEffectiveContinuity.lean} through the stronger
one-term recurrence route.\\
\addlinespace
Effective uniform continuity & Closed modulus and recursive integer arithmetic &
Proved in \cref{thm:reciprocal-modulus}.  The repository predicate itself is
formalized as \path{Fabius.fabiusInv_effectivelyUniformContinuous}, using the
simple primitive-recursive factorial denominator at $r=n$.  The exact smaller
logarithmic-$r(n)$ packaging is formalized by
\path{Fabius.inverseFabiusLogarithmicOrder_isLeast},
\path{Fabius.inverseFabiusLogarithmicDeltaDenominator_primrec}, and
\path{Fabius.fabiusInv_effectivelyUniformContinuous_logarithmicDelta}; its
stronger factorial companion is also formalized.\\
\addlinespace
Exact dyadic ceiling denominator $d_*$ & Exact rational dyadic evaluation,
ceiling arithmetic, and the exact threshold & Human proof in
\cref{prop:exact-dyadic-modulus}; neither the ceiling construction nor its
qualified fixed-target minimality is yet formalized.\\
\addlinespace
Sequential computability & Forward rational oracle, closed inverse modulus,
and tolerant bisection & Proved in
\cref{thm:bisection-correct,cor:GT-sequential}.  No tolerant-bisection or
sequential-computability declaration is supplied by the current inverse
effective-continuity modules.\\
\addlinespace
Optimal bit asymptotic & Closed lower bound plus iterated-flatness upper bound &
Paper-level proof in \cref{thm:mu-asymptotic}; no direct public Lean theorem.\\
\addlinespace
Sharp periodic refinement & Frontier all-orders endpoint inverse theory &
Imported as a paper-level input and reinterpreted in
\cref{cor:sharp-asymptotic-modulus}; the leading equivalent is formalized,
but the explicit periodic correction is not.\\
\bottomrule
\end{longtable}

This table is also a minimal trust boundary.  The main computability theorem
does not rely on the conjectures in \cref{sec:frontiers}, on numerical data,
or on any asymptotic expansion.

\section{Corpus crosswalk}
\label{app:corpus-crosswalk}

\begin{longtable}{@{}>{\raggedright\arraybackslash}p{0.30\textwidth}>{\raggedright\arraybackslash}p{0.61\textwidth}@{}}
\toprule
Corpus component & Role in the present report \\
\midrule
\endhead
\path{Fabius_Function_and_Rvachev_Up.tex} and archived primary expositions &
Normalization, probability law, up-function relation, exact dyadic arithmetic,
finite splines, and endpoint asymptotic context.\\
\addlinespace
\path{fabius_lean_walkthrough.tex} & Public theorem map; distinction among the
bounded CDF, signed extension, and up-function; forward computability and
inverse module architecture.\\
\addlinespace
\path{FabiusFunction_Mathematical_Glossary.tex} & Repository-specific
Grzegorczyk definitions and terminology for fast names, sequential
computability, and effective uniform continuity.\\
\addlinespace
\path{PAPER_COVERAGE.md} and the two Arias de Reyna paper transcriptions &
Formal/source crosswalk for probability, Fourier, dyadic, moment, and
arithmetic results.\\
\addlinespace
\path{FabiusComputability.lean} and
\path{FabiusComputableSpline.lean} & Existing forward $2$-Lipschitz theorem,
uniform centered-spline certificate, primitive-recursive evaluator, and
computable-real-function packaging.\\
\addlinespace
\path{FabiusInverse.lean} & Existing clamped inverse, inverse identities,
continuity, interior derivatives, curvature, and endpoint steepness.\\
\addlinespace
\path{Convexity.lean} and \path{Monotonicity.lean} & Formal strict density
shape and strict increase used as the geometric input to the paper-level
least-mass argument.\\
\addlinespace
\path{FabiusInverseAsymptotic.lean} & Formal leading endpoint equivalent and
implicit expansion pulled back to exact inverse coordinates; not an explicit
periodic or all-orders inverse reversion.\\
\addlinespace
\path{InverseModulus.lean} & Post-snapshot formalization of fixed-length
Fabius increments, their global weak and maximal strict shapes, the least
endpoint increment and complete equality locus, constrained forward
superadditivity with equality cases, local and global inverse-gap inequalities,
unit-input inverse-gap rigidity, global inverse subadditivity, absolute and
metric pointwise moduli, attained exact unit-interval and totalized suprema, and
the exact effective-injectivity threshold.  Exact declaration names appear in
\cref{sec:lean-status}.\\
\addlinespace
\path{FabiusInverseEffectiveContinuity.lean} & Post-snapshot formalization of
the recurrence-derived factorial endpoint lower bound, the exact numerical
$\Delta_r$ endpoint inequality, strict and closed versions of the Delta and
factorial dyadic inverse moduli, primitive-recursive denominator arithmetic,
and effective uniform continuity with the simple $r=n$ factorial witness.  It
does not formalize the probabilistic event, logarithmic $r(n)$, exact ceiling
$d_*$, or a sequential realizer.  Every public declaration is crosswalked in
\cref{sec:lean-status}.\\
\addlinespace
\path{FabiusInverseLogarithmicModulus.lean} & Post-snapshot formalization of
the primitive-recursive least order $r(n)$, its binary-length identity, the
report-exact $d(n)$ and stronger factorial denominator, strict and closed-input
reciprocal moduli, and effective uniform continuity with either logarithmic
witness.  It does not formalize the exact ceiling $d_*$, tolerant bisection, a
sequential realizer, or input-bit asymptotics.  Every public declaration is
crosswalked in \cref{sec:lean-status}.\\
\addlinespace
\path{Inverse_and_Sampling_Frontiers.tex} and
\path{Inverse_Endpoint_All_Orders.tex} & Exact finite-prefix quantiles,
$q$-Richardson filters, Bell-polynomial inversion, and explicit all-orders
endpoint quantile asymptotics at frontier-document, not Lean, status.\\
\addlinespace
\path{Small_Argument_Asymptotics.tex} & Explicit lower-Lambert forward
asymptotics and periodic Gamma--zeta coefficient calculus.\\
\addlinespace
$q$-calculus, sinc-product, Fourier-decay, Legendre/Lagrange, total-positivity,
and frontier-compilation sources & Context for alternative certified oracles,
shape strengthening, and research directions; not logical prerequisites of
the computability proof.\\
\bottomrule
\end{longtable}

The audit included generated \texttt{.tex} fragments and frozen variants, but
the report cites parent documents rather than treating generated rows as
independent mathematical sources.  The arrival hashes preserve the delivered
commit-\arrivalcommit\ view; the counts and formalization boundaries stated in
this filed edition come from the live-corpus intake rerun described above.

\section{Reproducing the exact-rational supplement}
\label{app:reproduction}

The archive accompanying this report contains
\begin{center}
\path{inverse_fabius_computability_experiments.py}
\end{center}
and its captured output \path{numerical_output.txt}.  The program uses only
Python's standard library.  From the extracted package directory, run

\begin{lstlisting}[style=reportpseudo,language={}]
python inverse_fabius_computability_experiments.py
\end{lstlisting}

The program performs all spline evaluations with
\texttt{fractions.Fraction}.  Its principal checks are:

\begin{enumerate}
\item the endpoint event lower bound and flatness upper bound in
      \cref{tab:endpoint-bounds};
\item exact convergence of the centered splines at $1/2$, $1/4$, and $1/8$,
      using the exact values $1/2$, $5/72$, and $1/288$;
\item a finite-grid analogue of the least-interval-mass theorem;
\item a short certified bisection transcript illustrating all comparison
      margins.
\end{enumerate}

These computations are regression tests and illustrations.  No numerical
observation is used as a premise of a theorem.

\begin{thebibliography}{99}

\bibitem{G1957}
A.~Grzegorczyk,
\newblock On the definitions of computable real continuous functions,
\newblock \emph{Fundamenta Mathematicae} \textbf{44} (1957), 61--71.

\bibitem{Weihrauch2000}
K.~Weihrauch,
\newblock \emph{Computable Analysis: An Introduction},
\newblock Springer, Berlin, 2000.

\bibitem{PourElRichards1989}
M.~B.~Pour-El and J.~I.~Richards,
\newblock \emph{Computability in Analysis and Physics},
\newblock Springer, Berlin, 1989.

\bibitem{Fabius1966}
J.~Fabius,
\newblock A probabilistic example of a nowhere analytic $C^\infty$-function,
\newblock \emph{Zeitschrift f\"ur Wahrscheinlichkeitstheorie und Verwandte
Gebiete} \textbf{5} (1966), no.~2, 173--174,
\newblock \url{https://doi.org/10.1007/BF00536652}.

\bibitem{JessenWintner1935}
B.~Jessen and A.~Wintner,
\newblock Distribution functions and the Riemann zeta function,
\newblock \emph{Transactions of the American Mathematical Society}
\textbf{38} (1935), no.~1, 48--88,
\newblock \url{https://doi.org/10.1090/S0002-9947-1935-1501802-5}.

\bibitem{Arias05442}
J.~Arias de Reyna,
\newblock An infinitely differentiable function with compact support:
definition and properties,
\newblock English translation of the 1982 paper,
\newblock \url{https://arxiv.org/abs/1702.05442}.

\bibitem{Arias06487}
J.~Arias de Reyna,
\newblock Arithmetic of the Fabius function,
\newblock \url{https://arxiv.org/abs/1702.06487}.

\bibitem{ProveIt}
V.~Reshetnikov,
\newblock \texttt{ProveIt}, Fabius-function development and documentation,
\newblock arrival repository snapshot \arrivalcommit,
\newblock \url{https://github.com/VladimirReshetnikov/ProveIt}.

\bibitem{WikipediaCRF}
Wikipedia contributors,
\newblock Computable real function,
\newblock accessed 30 August 2026,
\newblock \url{https://en.wikipedia.org/wiki/Computable_real_function}.

\end{thebibliography}

\end{document}
